\section*{四个公理}
\subsection*{公理2与个数问题}
\begin{enumerate}
    \item 过一条直线和这条直线外不共线的三点,最多可确定\dotfill{(\kh{B})}
    \xx{三个平面}{四个平面}{五个平面}{六个平面}
    \begin{jieda}
        记直线为 $l,l$ 外不共线的 3 点为 $P,Q,R$ . 根据公理 2 及其推论 1,直线 $l$ 和 $P,Q,R$ 任意一点各能确定 1 个平面; $P,Q,R$ 三点又能确定 1 个平面. 这 4 个平面可以两两不重合 (例: 正方体 ${ABCD} - {A}^{\prime }{B}^{\prime }{C}^{\prime }{D}^{\prime }$ 中, $D{D}^{\prime }$ 与 $A,B,C$ 三点),故最多可确定 4 个平面. 
        
        \ps 过一条直线和这条直线外不共线的三点,由于平面 ${PQR}$ 必然存在,因此最少可以确定 1 个平面,可以确定 1 或 3 或 4 个平面. 确定 1 个平面的情况是直线 $l$ 恰好在平面 ${PQR}$ 上且不过 $P,Q,R$ 任意一点. 确定 3 个平面的情况如: 正方体 ${ABCD} - {A}^{\prime }{B}^{\prime }{C}^{\prime }{D}^{\prime }$ 中, ${A}^{\prime }{B}^{\prime }$ 与 $A,B,C$ 三点.
    \end{jieda}
    \item 若三条直线两两相交, 则最多可以确定平面的个数是 \dotfill{(\kh{C})}
    \xx{1}{2}{3}{4}
    \begin{jieda}
        根据公理 2 的推论 2,任意两条相交直线可以确定一个平面,三条直线最多可以确定 3 个平面,且可以两两不重合(例:正方体 ${ABCD} - {A}^{\prime }{B}^{\prime }{C}^{\prime }{D}^{\prime }$ 中,直线 ${AB}$ 、直线 ${AD}$ 、直线 $A{A}^{\prime }$ ,确定了平面 ${ABCD}$ 、平面 ${AB}{B}^{\prime }{A}^{\prime }$ 、平面 ${AD}{D}^{\prime }{A}^{\prime }$ ). 
        
        \ps 三条直线 $l,m,n$ 两两相交,可以确定 1 或 3 个平面. 确定 1 个平面的情况是这 3 条直线共面.
    \end{jieda}
    
    \item \begin{enumerate}
        \item 若共面的三条直线两两相交,则交点的个数是\blkda{1 或 3}
        \begin{jieda}
            交点最少为 1 个,此时三条直线交于同一点: 最多为个 ${C}_{3}^{2} = 3$ 个,例如三角形三边所在直线. 仅有 2 个交点是不可能的,因为根据 “过两点有且只有一条直线” ,这必会导致这三条直线中的两条重合.
        \end{jieda}
        \item 四条直线过同一点,过每两条直线作一个平面,则可以作\blkda[2]{1 或 4 或 6}个不同的平面
        \begin{jieda}
            记直线为 $a,b,c,d$ ,根据共面直线条数,分以下 3 种情况讨论: 
            \begin{dingautolist}{172}
                \item 四条直线共面,只能作 1 个平面
                \item 三条直线共面,不失一般性,假设 $a,b,c \subset  {\alpha }$ 但 $d \not \subset {\alpha }$ . 由于直线 $a,b,c$ 互不重合,根据公理 3,直线 $d$ 与直线 $a$ 或 $b$ 或 $c$ 确定的平面也互不重合,且异于平面 $\alpha$ (例:正方体 ${ABCD} - {A}^{\prime }{B}^{\prime }{C}^{\prime }{D}^{\prime }$ 中,直线 ${AB}$ 、直线 ${AC}$ 、直线 ${AD}$ 、直线 $\left. {{{AA}}}^{\prime }\right)$ ,此时可作 4 个不同平面
                \item 任意三条直线都不共面,根据公理 2 的推论 2,此时最多能作 ${C}_{4}^{2} = 6$ 个平面 (例: 四棱锥的四条棱), 且这 6 个面两两不重合 (否则至少有三条直线共面), 因此可作 6 个不同平面.
            \end{dingautolist}
        \end{jieda}
        \item 互相平行的四条直线,每两条确定一个平面,最多可确定\blkda{6}个平面,最少可确定\blkda{1}个平面
        \begin{jieda}
            最多可确定 ${C}_{4}^{2} = 6$ 个平面 (例:正方体 ${ABCD} - {A}^{\prime }{B}^{\prime }{C}^{\prime }{D}^{\prime }$ 中,直线 $A{A}^{\prime }$ 、直线 $B{B}^{\prime }$ 、直线 $C{C}^{\prime }$ 、直线 $D{D}^{\prime }$ ). 最少可确定 1 个平面,此时四条直线在同一平面上. 
            
            \ps 仿照第(2)问,根据共面直线条数分类,可知互相平行的四条直线, 可以确定 1 个或 4 个或 6 个平面.
        \end{jieda}
    \end{enumerate}
\end{enumerate}
\subsection*{三点共线与三线共点}

\subsubsection*{三点共线}
\begin{enumerate}
    \item \begin{enumerate}
        \item 如图所示，三角形$ABC$在平面$\alpha$外，它的三边所在直线分别交平面$\alpha$于$P,Q,R$三点.求证：$P,Q,R$三点在同一条直线上.
        \begin{center}
            \includegraphics[width=0.25\linewidth]{bodyxb//src/sandiangongxian1.png}
        \end{center}
        \begin{jieda}
            $\left . \begin{array}{l}
                 P \in \text{面}ABC, P \in \alpha  \\
                 Q \in \text{面}ABC, Q \in \alpha \\
                 R \in \text{面}ABC, R \in \alpha
            \end{array} \right\} \Longrightarrow P,Q,R \in \alpha \cap \text{面}ABC$，故三点共线.
        \end{jieda}
        \item 在四边形 ${ABCD}$ 中,已知 ${AB}//{CD},{AB},{BC},{AD},{DC}$ (或延长线) 分别与平面 $\alpha$ 相交于点 $E,G,H,F$ (如图). 求证: 点 $E,F,G,H$ 必在同一条直线上
        \begin{center}
        \includegraphics[width=0.2\textwidth]{bodyxb/src/bo_d20aucref24c73anc48g_131_1138_1171_245_215_0.jpg}
        \end{center}
        \begin{jieda}
            证明: $\because {AB}//{CD},\therefore {AB},{CD}$ 确定一个平面 (公理 2 推论 3),记为 $\beta$ . 
            
            在平面 ${ABCD}$ 中, ${AB},{BC},{CD},{DA} \subset  \beta$ (公理 1), $\therefore E \in  {AB} \subset  \beta$ ,又 $E \in  \alpha ,\therefore E$ 在 $\alpha$ 与 $\beta$ 的交线上 (公理 3),记交线为 $l$ . 
            
            同理, $F,G,H \in  l,\therefore E,F,G,H$ 四点共线.
        \end{jieda}
    \end{enumerate}

    \item 如图，在三棱锥 $A - {BCD}$ 中，作截面 ${PQR}$ ， ${PQ}$ ， ${CB}$ 的延长线交于点 $M$ ， ${RQ}$ ， ${DB}$ 的延长线交于点 $N,{RP},{DC}$ 的延长线交于点 $K$ . 判断 $M,N,K$ 三点是否共线，并说明理由.

    \begin{center}
    \includegraphics[width=0.2\textwidth]{bodyxb/src/bo_d4k1jdbef24c739co0cg_0_244_601_360_357_0.jpg}
    \end{center}
    \begin{jieda}
        $M,N,K$ 三点共线. 理由如下:

        因为 $M\text{ 、 }N$ 即在平面 ${BCD}$ 内又在平面 ${PRQ}$ 内,所以 $M\text{ 、 }N$ 在平面 ${BCD}$ 与平面 ${PRQ}$ 的交线上,所以 ${MN}$ 是平面 ${BCD}$ 与平面 ${PRQ}$ 的交线,
        
        $N\text{ 、 }K$ 即在平面 ${BCD}$ 内又在平面 ${NKR}$ 内,所以 $N\text{ 、 }K$ 在平面 ${BCD}$ 与平面 ${NKR}$ 的交线上,所以 ${NK}$ 是平面 ${BCD}$ 与平面 ${NKR}$ 的交线,
        
        又平面 ${NKR}$ 与平面 ${PRQ}$ 是同一平面,所以 ${MN}$ 与 ${NK}$ 是同一条直线,即 $M,N,K$ 三点共线.
    \end{jieda}
    \item 如图所示,在正方体 ${ABCD} - {A}_{1}{B}_{1}{C}_{1}{D}_{1}$ 中,点 $O$ 为正方形 ${ABCD}$ 的中心, $H$ 为直线 ${B}_{1}D$ 与平面 ${AC}{D}_{1}$ 的交点. 求证: ${D}_{1}\text{ 、 }H\text{ 、 }O$ 三点共线.
    
    \begin{center}
    \includegraphics[width=0.2\textwidth]{bodyxb/src/bo_d1i8hsn7aajc7382r7dg_5_406_1683_258_275_0.jpg}
    \end{center}
    \begin{jieda}
        连接 ${BD}$ 、 ${B}_{1}{D}_{1}$

        因为${B}_{1}D \cap \text{平面} {AC}{D}_{1} = H$，所以$H \in $平面$BB_1D_1D$且$H \in$平面 ${AC}{D}_{1}$
        
        因为点 $O$ 为正方形 ${ABCD}$ 的中心，所以$O \in BD, O \in AC$，因此$O \in $平面$BB_1D_1D$且$O \in$平面 ${AC}{D}_{1}$

        易知所以$D_1 \in$平面$BB_1D_1D$且$D_1 \in$平面 ${AC}{D}_{1}$

        根据公理3可知三点共线
    \end{jieda}
\end{enumerate}
\subsubsection*{三线共点}
三线共点问题的证明思路：首先根据两条相交直线确定交点，再证明交点在另一条直线上.

\begin{enumerate}
    \item 如图所示，已知平面$\alpha, \beta$，且有$\alpha \cap \beta = l$，设在梯形$ABCD$中，$AD // BC$，且$AB \subset \alpha, CD \subset \beta$，求证：$AB, CD, l$共点.

    \begin{center}
    \includegraphics[width=0.2\textwidth]{bodyxb/src/bo_d4k1oc3ef24c739co680_0_243_599_371_239_0.jpg}
    \end{center}
    \begin{jieda}
        如图,梯形 ${ABCD}$ 中，因为 ${AD}//{BC}$，所以 ${AB}$ 与 ${CD}$ 必交于一点,
        
        设 ${AB}$ 交 ${CD}$ 于点 $M$ ,则 $M \in  {AB},M \in  {CD}$ ,
        
        又因为 ${AB} \subset  \alpha ,{CD} \subset  \beta$，所以 $M \in  \alpha ,M \in  \beta$ ,
        
        又因为 $\alpha  \cap  \beta  = l$ ,所以 $M \in  l$，所以 ${AB},{CD},l$ 共点.
        
        \begin{center}
        \includegraphics[width=0.2\textwidth]{bodyxb/src/bo_d4k1oc3ef24c739co680_0_240_1495_379_264_0.jpg}
        \end{center}
    \end{jieda}
    
    \item 如图， $\bigtriangleup  {{ABC}}$ 和 $\bigtriangleup  {{A}}_{1}{{B}}_{1}{{C}}_{1}$ 不在同一个平面内，如果三直线 ${A}{{A}}_{1}\text{ 、 }{B}{{B}}_{1}\text{ 、 }{{{CC}}}_{1}$ 两两相交，求证: 三直线 $A{A}_{1}\text{ 、 }B{B}_{1}\text{ 、 }C{C}_{1}$ 交于一点.

    \begin{center}
    \includegraphics[width=0.2\textwidth]{bodyxb/src/bo_d4g0jos601uc73fkjrl0_3_246_337_358_238_0.jpg}
    \end{center}

    \begin{jieda}
        设 $B{B}_{1}$ 与 $C{C}_{1},C{C}_{1}$ 与 $A{A}_{1},A{A}_{1}$ 与 $B{B}_{1}$ 分别确定平面 $\alpha ,\beta ,\gamma$，$A{A}_{1}$ 与 $B{B}_{1}$ 的交点为 $P$
        
        因为 $P \in  A{A}_{1},P \in  B{B}_{1}$，$A{A}_{1} \subset  \beta ,B{B}_{1} \subset  \alpha$ ,所以 $P \in  \alpha ,P \in  \beta$ ,即 $P \in  \alpha  \cap  \beta$ .
        
        又 $\alpha  \cap  \beta  = C{C}_{1}$ ,所以 $P \in  C{C}_{1}$ .所以三条直线 $A{A}_{1},B{B}_{1},C{C}_{1}$ 交于一点
    \end{jieda}
    
    \item 如图所示，在空间四边形 ${ABCD}$ 中, $H\text{ 、 }G$ 分别是 ${AD}$ 、 ${CD}$ 的中点, $E\text{ 、 }F$ 分别是边 ${AB}\text{ 、 }{BC}$ 上的点,且 $\dfrac{CF}{FB} = \dfrac{AE}{EB} = \dfrac{1}{3}$ . 求证: 直线 ${EH}\text{ 、 }{BD}\text{ 、 }{FG}$ 相交于一点.
    
    \begin{center}
    \includegraphics[width=0.2\textwidth]{bodyxb/src/bo_d1i8hsn7aajc7382r7dg_11_1209_854_284_272_0.jpg}
    \end{center}

    \begin{jieda}
        证明 
        
        如图所示，连接 ${EF}\text{ 、 }{GH}$ .
        
        在三角形 ${ACD}$ 中,因为 $H\text{ 、 }G$ 分别是 ${AD}\text{ 、 }{CD}$ 的中点,所以 ${GH}//{AC}$ ,且 ${GH} = \dfrac{1}{2}{AC}$ .
        
        \begin{center}
        \includegraphics[width=0.2\textwidth]{bodyxb/src/bo_d1i8hsn7aajc7382r7dg_11_1112_1181_382_267_0.jpg}
        \end{center}
        
        在三角形 ${ABC}$ 中,因为 $\dfrac{CF}{FB} = \dfrac{AE}{EB} = \dfrac{1}{3}$ ,所以 ${EF}//{AC}$ ,且 ${EF} = \dfrac{3}{4}{AC}$ .
        
        由公理 4,得 ${GH}//{EF}$ ,所以 $E\text{ 、 }F\text{ 、 }G\text{ 、 }H$ 四点共面.
        
        又因为 ${GH} \neq  {EF}$ ,所以 ${EH}$ 与 ${FG}$ 相交,设交点为 $P$ . 因为 ${EH} \subset$ 平面 ${ABD}$ ,所以 $P \in$ 平面 ${ABD}$ ,同理 $P \in$ 平面 ${BCD}$ . 又由平面 ${ABD} \cap$ 平面 ${BCD} = {BD}$ ,因此 $P \in  {BD}$ . 所以,直线 ${EH}\text{ 、 }{BD}\text{ 、 }{FG}$ 相交于一点.
    \end{jieda}
\end{enumerate}

\section*{线线位置关系}
\subsection*{异面直线}
\begin{enumerate}
    \item 两条异面直线是指\dotfill{(\kh{D})}
    \xx{空间两条没有公共点的直线}{平面上一直线与这平面外的一条直线}{分别在两个平面上的两条直线}{不同在任何一个平面上的两条直线}
    \begin{jieda}
        选项 (D) 是异面直线的定义. 其余 3 个选项的反例如下:对于选项 (A),两条平行直线；对于选项 (B),正方体 ${ABCD} - {A}^{\prime }{B}^{\prime }{C}^{\prime }{D}^{\prime }$ 中,平面 ${ABCD}$ 内的直线 ${AB}$ 与平面外的直线 ${A}^{\prime }{B}^{\prime }$ ; 对于选项 (D),正方体 ${ABCD} - {A}^{\prime }{B}^{\prime }{C}^{\prime }{D}^{\prime }$ 中,平面 ${ABCD}$ 内的直线 ${AB}$ 与平面 ${A}^{\prime }{B}^{\prime }{C}^{\prime }{D}^{\prime }$ 内的直线 ${A}^{\prime }{B}^{\prime }$
    \end{jieda}
    \item 已知平面 $\alpha  \cap$ 平面 $\beta  =$ 直线 $c$ ,直线 $a \subset  \alpha ,a//c$ ,直线 $b \subset  \beta$ ,且 $b$ 与 $c$ 相交,则 $a$ 和 $b$ 的位置关系是\dotfill{(\kh{C})}
    \xx{平行}{相交}{异面}{上述三种情况都有可能}
    \begin{jieda}
        符合异面直线判定定理的条件
    \end{jieda}
    \item 已知 $a,b$ 是异面直线,直线 $c//a$ ,那么 $c$ 与 $b$ \dotfill{(\kh{C})}
    \xx{一定是异面直线}{一定是相交直线}{不可能是平行直线}{不可能是相交直线}
    \begin{jieda}
        根据公理 4 ; 空间内,平行于同一直线的两直线互相平行. $c//b$ 是不可能的. 直线 $c$ 和 $b$ 异面和相交均有可能,例如,正方体 ${ABCD} - {A}^{\prime }{B}^{\prime }{C}^{\prime }{D}^{\prime }$ 中,取 $a = {AB},b = {A}_{1}{D}_{1},c = {CD}$ 时, $b$ 与 $c$ 异面; $c = {D}_{1}{C}_{1}$ 时, $b$ 与 $c$ 相交. 因此只有选项 (C) 正确.
    \end{jieda}
    \item 若 $a,b$ 是异面直线, $b,c$ 也是异面直线,则 $a,c$ 的位置关系是\dotfill{(\kh{A})}
    \xx{相交、平行或异面}{相交或平行}{异面}{平行或异面}
    \begin{jieda}
        以正方体 ${ABCD} - {A}^{\prime }{B}^{\prime }{C}^{\prime }{D}^{\prime }$ 为例,取 $a = {AB},b = {A}_{1}{D}_{1}$ ,若 $c = {BD}$ ,则 $a$ , $c$ 相交；若 $c = {CD}$ ,则 $a$ , $c$ 平行；若 $c =$  $C{C}_{1}$ ,则 $a,c$ 异面.
    \end{jieda}
    \item 经过空间一点 $A$ ,作与直线 $l$ 成 $\dfrac{\pi }{3}$ 角的直线共有 \dotfill{(\kh{D})}
    \xx{2 条}{3 条}{4 条}{无数条}
    \begin{jieda}
        若直线 $l$ 不经过点 $A$ ,则过 $A$ 点作直线 ${l}^{\prime }//l$ ,若直线 $l$ 经过点 $A$ ,取 ${l}^{\prime } = l$ . 在空间中,任作一条过点 $A$ 且与直线 ${l}^{\prime }$ 成 $\dfrac{\pi }{3}$ 角的直线,再将该直线绕着直线 ${l}^{\prime }$ 旋转一周 (得到两个对顶的圆锥),即得到无数条符合要求的直线.
    \end{jieda}
    \item 已知异面直线 $a$ 与 $b$ 满足 $a \subset  \alpha ,b \subset  \beta$ ,且 $\alpha  \cap  \beta  = l$ ,则直线 $l$ 与 $a,b$ 的位置关系一定是\dotfill{(\kh{B})}
    \xx{$l$ 与 $a,b$ 都相交}{$l$ 至少与 $a,b$ 中的一条相交}{$l$ 至多与 $a,b$ 中的一条相交}{$l$ 至少与 $a,b$ 中的一条平行}
    \begin{jieda}
        由于 $l \subset  \alpha ,l \subset  \beta$ ,若 $l$ 与直线 $a,b$ 均无交点,则 $a//l,b//l$ ,根据公理 $4,a//b$ ,这与 $a,b$ 是异面直线矛盾,故 $l$ 至少与 $a,b$ 中的一条相交,选项 (B) 正确,而选项 (C) 错误. 以正方体 ${ABCD} - {A}^{\prime }{B}^{\prime }{C}^{\prime }{D}^{\prime }$ 为例,取 $\alpha  =$ 平面 ${ABCD},\beta  =$ 平面 ${AD}{D}^{\prime }{A}^{\prime }$ ,则 $l = {AD}$ ,选项 (A) 反例: 取 $a = {AB},b = {A}^{\prime }{D}^{\prime }$ ; 选项 (D) 的反例: 取 $a = {AB},b = {A}^{\prime }D$
    \end{jieda}
    \item 有以下三个命题:\ding{172}若两条直线和第三条直线成等角,则这两条直线互相平行；\ding{173}若两条直线都和第三条直线垂直,则这两条直线互相平行；\ding{174}若两条直线都和第三条直线平行, 则这两条直线互相平行. 其中不正确的个数为\dotfill{(\kh{C})}
    \xx{0}{1}{2}{3}
    \begin{jieda}
        \ding{174}是公理4，\ding{172}\ding{173}均有异面的可能
    \end{jieda}
    \item 分别和两条异面直线都相交的两直线一定是\dotfill{(\kh{A})}
    \xx{不平行的直线}{不相交的直线}{相交直线或平行直线}{既不相交又不平行的直线}
    \begin{jieda}
        如果和两条异面直线都相交的两直线平行, 由于两平行直线可确定一个平面, 根据公理 1, 推知这两条异面直线在同一平面内,矛盾,故选项 $A$ 正确,选项 $C$ 错误. 以正方体 ${ABCD} - {A}^{\prime }{B}^{\prime }{C}^{\prime }{D}^{\prime }$ 为例, ${AB}$ 和 ${A}^{\prime }{D}^{\prime }$ 是两条异面直线,与它们都相交的两直线 $A{A}^{\prime },A{D}^{\prime }$ 相交,而 $A{A}^{\prime },B{D}^{\prime }$ 异面,故选项 $B,D$ 均错误.
    \end{jieda}
    \item 在正方体 ${ABCD} - {A}_{1}{B}_{1}{C}_{1}{D}_{1}$ 中,各面对角线所在的直线中与 ${B}_{1}D$ 成异面直线的条数是\dotfill{(\kh{D})}
    \xx{3}{4}{5}{6}
    \begin{jieda}
        各个面总计 12 条对角线,首先排除从点 ${B}_{1}$ 或点 $D$ 出发的 6 条,其余 6 条可使用异面直线判定定理 (平面外一点与平面内一点的连线,和平面内不经过该点的直线是异面直线),观察对角线、对角线所在的侧面和 ${B}_{1}D$ ,不难判定剩余 6 条都与 ${B}_{1}D$ 异面.
    \end{jieda}
\end{enumerate}
\subsection*{线线角}
\begin{enumerate}
    \item 在正方体 ${ABCD} - {A}_{1}{B}_{1}{C}_{1}{D}_{1}$ 中,
    \begin{enumerate}
        \item 各个面的对角线所在直线中,与 $A{B}_{1}$ 异面且所成角为 ${60}^{ \circ  }$ 的直线有\blkda{4}条
        \begin{jieda}
            各个面总计 12 条对角线,与 $A{B}_{1}$ 成异面直线的是 $B{C}_{1},{A}_{1}D,{A}_{1}{C}_{1},{BD},{D}_{1}C$ ,通过平移,可知前 4 条与 $A{B}_{1}$ 成 ${60}^{ \circ  }$ 角,而 ${D}_{1}C$ 与 $A{B}_{1}$ 成 ${45}^{ \circ  }$ 角,因此共 4 条满足题意.
        \end{jieda}
        \item 若 $E$ , $F$ 分别为棱 ${A}_{1}{B}_{1}$ , $B{B}_{1}$ 的中点,则 ${AE}$ 和 ${CF}$ 所成角的余弦值为\blkda{$\dfrac{2}{5}$}
        \begin{jieda}
            取平面$CC_1D_1D$的中点$O$，连接$EO$，$EO // CF$，则$\angle{AEO}$即为异面直线所成角（或其补角），根据余弦定理可知$cos\angle{AEO} = \dfrac{5+5-6}{10} = \dfrac{2}{5}$
        \end{jieda}
        \item 若 $G$ 为 ${CD}$ 的中点,则异面直线 ${B}_{1}C$ 与 ${AG}$ 所成角的正弦值是\blkda{$\dfrac{\sqrt{15}}{5}$}
        \begin{jieda}
            取$AB$中点$H$，连接$CH$，$CH // AG$，故$\angle{B_1CH}$即为异面直线所成角（或其补角），根据余弦定理可知$cos\angle{B_1CH} = \dfrac{5+8-5}{4\sqrt{10}} = \dfrac{\sqrt{10}}{5}$，故$sin\angle{B_1CH} = \dfrac{\sqrt{15}}{5}$
        \end{jieda}
    \end{enumerate}
    \item 在正方体 ${ABCD} - {A}_{1}{B}_{1}{C}_{1}{D}_{1}$ 中,
    \begin{enumerate}
        \item 若 $F$ , $G$ 分别是棱 $B{B}_{1}$ , ${DC}$ 的中点,则 ${AF}$ 与 ${D}_{1}G$ 所成角为\blkda{${90}^{ \circ  }$}
        \begin{jieda}
            取$CC_1$中点$H$，$DH // AF$，设$D_1G$与$DH$交于$M$，则$\angle{DMD_1}$即为异面直线所成角（或其补角），根据勾股定理可知$\angle{DMD_1} = 90^{\circ}$
        \end{jieda}
        \item 已知 $E,F$ 分别是棱 $B{B}_{1},C{C}_{1}$ 的中点,则异面直线 ${AE},{BF}$ 所成角的余弦值等于\blkda{$\dfrac{1}{5}$}
        \begin{jieda}
            取$DD_1$中点$H$，$AH // BF$，则$\angle{HAE}$即为异面直线所成角（或其补角），根据余弦定理可知$cos\angle{HAE} = \dfrac{5+5-8}{10} = \dfrac{1}{5}$
        \end{jieda}
    \end{enumerate}
    \item  在长方体 ${ABCD} - {A}_{1}{B}_{1}{C}_{1}{D}_{1}$ 中,
    
    \begin{center}
    \includegraphics[width=0.15\textwidth]{bodyxb/src/bo_d20aucref24c73anc48g_134_1256_820_186_303_0.jpg}
    \end{center}
    \begin{enumerate}
        \item 若棱 $B{B}_{1} = {BC} = 1,{AB} = \sqrt{3}$ ,则 $A{D}_{1}$ 与 ${BC}$ 所成角等于\blkda{${45}^{ \circ  }$}, $C{D}_{1}$ 与 ${AB}$ 所成角等于\blkda{${30}^{ \circ  }$}, $C{D}_{1}$ 与 ${A}_{1}D$ 所成角的正弦值等于\blkda{$\dfrac{\sqrt{14}}{4}$}
        \begin{jieda}
            ${BC} \bot  {AD}$ ,故 $A{D}_{1}$ 与 ${BC}$ 所成角为 $\angle {D}_{1}{AD} = \arctan \dfrac{{D}_{1}D}{AD} = \arctan \dfrac{1}{1} = {45}^{ \circ  }$
            
            ${AB} \bot  {CD}$ ,故 $C{D}_{1}$ 与 ${AB}$ 所成角为 $\angle {D}_{1}{CD} = \arctan \dfrac{{D}_{1}D}{CD} = \arctan \dfrac{1}{\sqrt{3}} = {30}^{ \circ  }$
            
            易证四边形 ${A}_{1}{B}_{1}{CD}$ 是平行四边形,故 ${A}_{1}D\bot {B}_{1}C$，不难求得 ${B}_{1}C = \sqrt{2},{D}_{1}C = {D}_{1}{B}_{1} = 2,C{D}_{1}$ 与 ${A}_{1}D$ 所成角为 $\angle {D}_{1}C{B}_{1},\cos \angle {D}_{1}C{B}_{1} = \dfrac{C{D}_{1}^{2} + C{B}_{1}^{2} - {B}_{1}{D}_{1}^{2}}{2 \cdot  C{D}_{1} \cdot  C{B}_{1}} = \dfrac{\sqrt{2}}{4}$ ,故 $\sin \angle {D}_{1}{CD} = \dfrac{\sqrt{14}}{4}.$
        \end{jieda}
        \item 若 $\angle {BA}{B}_{1} = \angle {B}_{1}{A}_{1}{C}_{1} = {60}^{ \circ  }$，则 $A{A}_{1}$ 与 ${B}_{1}C$ 所成角的余弦值等于\blkda{$\dfrac{\sqrt{2}}{2}$}, ${B}_{1}C$ 与 ${A}_{1}{C}_{1}$ 所成角的余弦值等于\blkda{$\dfrac{\sqrt{6}}{4}$}
        \begin{jieda}
            $\angle{BB_1C}$即为异面直线所成角（或其补角） $\angle B{B}_{1}C = {45}^{ \circ  },\cos \angle B{B}_{1}C = \dfrac{\sqrt{2}}{2}$ .
            
            $\angle{DA_1C_1}$即为异面直线所成角（或其补角） $\cos \angle DA_1C_1 = \dfrac{6+4-4}{4\sqrt{6}} = \dfrac{\sqrt{6}}{4}$ .
        \end{jieda}
        \item 若底面边长 ${AB} = {BC} = 3$ ,棱 $A{A}_{1} = 4$ ,则异面直线 $A{B}_{1}$ 与 ${A}_{1}D$ 所成角的余弦值等于\blkda{$\dfrac{16}{25}$}
        \begin{jieda}
            $\angle{AB_1C}$即为异面直线所成角（或其补角） $\cos \angle AB_1C = \dfrac{25+25-18}{50} = \dfrac{16}{25}$ 
        \end{jieda}
    \end{enumerate}
    \item 在四面体 ${ABCD}$ 中,
    \begin{enumerate}
        \item 若棱 ${AB} = {CD} = 2,E,F$ 分别是棱 ${AC},{BD}$ 的中点,且 ${EF} = \sqrt{3}$ ,则 ${AB}$ 与 ${CD}$ 所成角等于\blkda{$60^\circ$}
        \begin{jieda}
            取$BC$中点$G$，连接$EG, FG$，$\angle{EGF}$即为异面直线所成角（或其补角），$\cos \angle EGF = \dfrac{1+1-3}{2} = -\dfrac{1}{2}$，故异面直线所成角为$60^\circ$
        \end{jieda}
        \item 若 ${AC}$ 与 ${BD}$ 所成角为 ${60}^{ \circ  }$ ,且 ${AC} = {BD} = a$ ,则连接 ${AB},{BC},{CD},{DA}$ 这四条棱的中点所得四边形的面积等于\blkda{$\dfrac{\sqrt{3}}{8}{a}^{2}$}
        \begin{jieda}
            根据几何关系可知${AB},{BC},{CD},{DA}$ 这四条棱的中点所得四边形为边长为$\dfrac{a}{2}$的菱形，该菱形的的内角即异面直线所成角（或其补角），故其面积$S = \dfrac{1}{2}\times \dfrac{a}{2} \times \dfrac{\sqrt{3}a}{2} = \dfrac{\sqrt{3}}{8}{a}^{2}$
        \end{jieda}
    \end{enumerate}
    \item 正方体的两条体对角线相交, 其夹角的正弦值等于 \dotfill{(\kh{C})}
    \xx{$\dfrac{\sqrt{2}}{2}$}{$\dfrac{1}{3}$}{$\dfrac{2\sqrt{2}}{3}$}{$\dfrac{\sqrt{10}}{10}$}
    \begin{jieda}
        设正方体棱长为 1 . 正方体共有 4 条体对角线,任意两条都共面,且是长、宽分别为 $\sqrt{2}$ 、 1 的矩形的对角线. 对角线长 $\sqrt{3}$ ,且两条对角线分矩形为 4 个相等面积的三角形,设对角线所成角为 $\theta ,\therefore 4 \cdot  \dfrac{1}{2} \cdot  {\left( \dfrac{\sqrt{3}}{2}\right) }^{2} \cdot  \sin \theta  = \sqrt{2} \cdot  1$ ,
        $\therefore \sin \theta  = \dfrac{2\sqrt{2}}{3}$ (也可用余弦定理结合三角恒等式求解).
    \end{jieda}
    \item 在正方体 ${EFGH} - {E}^{\prime }{F}^{\prime }{G}^{\prime }{H}^{\prime }$ 中, $F{G}^{\prime }$ 与 ${EG}$ 所成的角等于 \dotfill{(\kh{B})}
    \xx{${45}^{ \circ  }$}{${60}^{ \circ  }$}{${90}^{ \circ  }$}{${120}^{ \circ  }$}
    \begin{jieda}
        易证 ${EF}{G}^{\prime }{H}^{\prime }$ 是平行四边形,故 $E{H}^{\prime }//F{G}^{\prime }$ . 而 $\bigtriangleup {EG}{H}^{\prime }$ 是等边三角形,故 $F{G}^{\prime }$ 与 ${EG}$ 所成的角为 $\angle {GE}{H}^{\prime } = {60}^{ \circ  }$ .
    \end{jieda}

    \item 在正四面体 ${ABCD}$ 中,
    \threetu{\includegraphics[width=0.6\textwidth]{bodyxb/src/bo_d20aucref24c73anc48g_136_277_839_266_240_0.jpg}}{\includegraphics[width=0.6\textwidth]{bodyxb/src/bo_d20aucref24c73anc48g_136_652_841_261_237_0.jpg}}{\includegraphics[width=0.6\textwidth]{bodyxb/src/bo_d20aucref24c73anc48g_136_1040_837_267_239_0.jpg}}
    \begin{enumerate}
        \item 若 $E$ 为棱 ${CD}$ 的中点 (如图),求 ${AE}$ 和 ${BC}$ 所成角的余弦值
        \begin{jieda}
            取$BD$中点$F$，则$\angle{AEF}$即为异面直线所成角（或其补角），根据余弦定理$cos\angle{AEF} = \dfrac{1+3-3}{2\sqrt{3}} = \dfrac{\sqrt{3}}{6}$
        \end{jieda}
        \item 已知 $E,F$ 分别是棱 ${CD},{AD}$ 的中点(如图),求 ${AE},{BF}$ 所成角的余弦值
        \begin{jieda}
            取$ED$中点$G$，则$\angle{BFG}$即为异面直线所成角（或其补角），根据余弦定理$cos\angle{BFG} = \dfrac{3+\dfrac{3}{4}-\dfrac{13}{4}}{3} = \dfrac{1}{6}$
        \end{jieda}
        \item 若 $E$ , $F$ 分别为棱 ${AB}$ , ${CD}$ 的中点(如图),求 ${AF}$ 与 ${CE}$ 所成角的正切值
        \begin{jieda}
            取$BF$中点$G$，则$\angle{CEG}$即为异面直线所成角（或其补角），根据余弦定理$cos\angle{BFG} = \dfrac{3+\dfrac{3}{4}-\dfrac{7}{4}}{3} = \dfrac{2}{3}$，故$tan\angle{BFG} = \dfrac{\sqrt{5}}{2}$
        \end{jieda}
    \end{enumerate}
\end{enumerate}

\section*{线面平行}
\subsection*{位置关系}
\begin{enumerate}
    
    \item 已知直线 $l//$ 平面 $\alpha$ ,直线 $a \subset  \alpha$ ,则 $l$ 与 $a$ 的位置关系必定是\dotfill{(\kh{D})}
    \xx{$l//a$}{$l$ 与 $a$ 异面}{$l$ 与 $a$ 相交}{$l$ 与 $a$ 无公共点}
    \begin{jieda}
        由线面平行的定义, $l \cap  \alpha  = \varnothing$ ,而 $a \subset  \alpha ,\therefore l \cap  a = \varnothing$ ,选项 $D$ 正确,选项 $C$ 错误. 在正方体 ${ABCD} -$  ${A}^{\prime }{B}^{\prime }{C}^{\prime }{D}^{\prime }$ 中,取 $l = {B}^{\prime }{C}^{\prime },\alpha  = {ABCD}$ ,若 $a = {BC}$ ,则 $l//a$ ; 若 $a = {AB}$ ,则 $l$ 与 $a$ 异面,故选项 $A,B$ 均错误.
    \end{jieda}
    \item $a,b$ 是两条直线, $\alpha$ 是平面. 给出以下四个命题:
    \begin{dingautolist}{172}
        \item 若 $a//\alpha ,b//\alpha$ ,则 $a//b$
        \item 若 $a//b,a//\alpha$ , 则 $b//\alpha$
        \item 若 $a//\alpha$ ,则 $a//\alpha$ 上所有直线
        \item 若 $a//\alpha$ 上无数条直线,则 $a//\alpha$
    \end{dingautolist}
    其中正确的个数是\dotfill{(\kh{A})}
    \xx{0}{1}{2}{3}
    \begin{jieda}
        \begin{dingautolist}{172}
            \item 三种位置关系均有可能
            \item $b$可能在$\alpha$内
            \item 存在异面的直线
            \item $a$可能在$\alpha$内
        \end{dingautolist}
    \end{jieda}
    \item 若直线 $a$ 平行于平面 $\alpha$ 上的直线 $b$ ,则 $a$ 和 $\alpha$ 的位置关系是\dotfill{(\kh{C})}
    \xx{$a//\alpha$}{$a \subset  \alpha$}{$a//\alpha$ 或 $a \subset  \alpha$}{$a//\alpha$ 或 $a$ 与 $\alpha$ 相交}
    \begin{jieda}
        若直线 $a$ 不在平面 $\alpha$ 上,则 $a//\alpha$ ; 若直线 $a$ 在平面 $\alpha$ 上,则 $a \subset  \alpha$
    \end{jieda}
    \item 已知不重合的直线 $a,b,c$ 及平面 $\alpha$ ,则下列条件能使 $a//b$ 成立的是 \dotfill{(\kh{C})}
    \xx{$a//\alpha$ 且 $b//\alpha$}{$a\bot c$ 且 $b\bot c$}{$a//c$ 且 $b//c$}{$a//\alpha$ 且 $a\bot c$}
    \begin{jieda}
        根据公理 4,选项 C 能推出 $a//b$ .
    \end{jieda}
    \item 给出四个命题:
    \begin{dingautolist}{172}
        \item ${AB}$ 为平面 $\alpha$ 外的线段,若点 $A,B$ 到平面 $\alpha$ 的距离相等,则 ${AB}//\alpha$
        \item 若一个角的两边分别平行于另一个角的两边,则这两个角相等
        \item 若直线 $a//$ 直线 $b$ ,则 $a$ 平行于过 $b$ 的所有平面
        \item 若直线 $a//$ 平面 $\alpha$ ,直线 $b//$ 平面 $\alpha$ ,则 $a//b$
    \end{dingautolist}
    其中正确的个数是\dotfill{(\kh{A})}
    \xx{0}{1}{2}{3}
    \begin{jieda}
        \begin{dingautolist}{172}
            \item 注意 $A,B$ 两点分别在平面 $\alpha$ 两侧的可能性
            \item 注意这两个角互补的可能性
            \item 直线 $a,b$ 本身可以确定一个平面,此时 $a$ 在该平面内
            \item 三种位置关系都有可能
        \end{dingautolist}
    \end{jieda}

    \item 如果点 $M$ 是两条异面直线 $a,b$ 外的一点,则过点 $M$ 且与 $a,b$ 都平行的平面 \dotfill{(\kh{C})}
    \xx{只有一个}{恰有两个}{或者没有, 或者只有一个}{有无数个}
    \item 已知 $A,B,C,D$ 是空间不共面的四点,它们到平面 $\alpha$ 的距离相等,则满足条件的平面 $\alpha$ 的个数是\dotfill{(\kh{C})}
    \xx{3}{4}{7}{8}
    \begin{jieda}
        按平面 $\alpha$ 两侧点的数量讨论: 
        \begin{dingautolist}{172}
            \item 若 4 点都在平面 $\alpha$ 的同侧,则四点共面,与题干矛盾
            \item 若只有 3 点在平面 $\alpha$ 的同侧有 4 种可能,现以同侧 3 点为点 $A,B,C$ 为例,此时平面 $\alpha //$ 平面 ${ABC}$ ,再加上 “距离相等” 的约束,故所求平面是唯一的,具体而言: 取线段 ${DA},{DB},{DC}$ 的中点,连接所得的平面即为所求
            \item 若只有 2 点在平面 $\alpha$ 的同侧,有 3 种可能,现以点 $A,B$ 和点 $C,D$ 分属平面 $\alpha$ 的两侧为例,此时 ${AB}//$ 平面 $\alpha ,{CD}//$ 平面 $\alpha ,{AB}$ 与 ${CD}$ 的公垂线段上平面 $\alpha$ ,“距离相等”的约束要求平面 $\alpha$ 经过 ${AB}$ 与 ${CD}$ 的公垂线段的中点
        \end{dingautolist}
        综上,满足条件的平面个数为 7 个.
    \end{jieda}

    \item \begin{enumerate}
        \item 若直线 $l$ 与平面 $\alpha$ 平行,则在平面 $\alpha$ 内与 $l$ 平行的直线有\blkda{无数}条
        \begin{jieda}
            过直线 $l$ 可以作无数个互不重合的平面,它们与平面 $\alpha$ 的交线 (如有) 有无数条,也互不重合,且都与直线 $l$ 平行；或者,过直线 $l$ 作 1 个平面与平面 $\alpha$ 相交,交线记为 ${l}^{\prime }$ ,而在平面 $\alpha$ 内有无数条直线与 ${l}^{\prime }$ 平行,根据公理 4, 它们都与直线 $l$ 平行.
        \end{jieda}
        \item 若 $P$ 是平面 $\alpha$ 外一点,则过 $P$ 作 $\alpha$ 的平行线有\blkda{无数}条
        \begin{jieda}
            在平面 $\alpha$ 内先取一点 ${P}^{\prime }$ ,再任取一点 ${Q}^{\prime }$ 异于点 ${P}^{\prime }$ ,则过 ${P}^{\prime },{Q}^{\prime },P$ 三点可确定一个平面,记为平面 $\beta$ . 在平面 $\beta$ 内,过点 $P$ 可以作一条直线 ${PQ}//{P}^{\prime }{Q}^{\prime }$ ,此时 ${PQ}//$ 平面 $\alpha$ . 当射线 ${P}^{\prime }{Q}^{\prime }$ 的方向不断变化时, ${PQ}$ 的方向也不断变化,便能得到无数条过点 $P$ 与平面 $\alpha$ 平行的直线. 或者,在学过两个平面平行的知识后,结合立体解析几何,可知: 过点 $P$ 可作唯一平面 $\beta$ 与平面 $\alpha$ 平行 (这两个平面的法向量相同),而平面 $\beta$ 中任意过点 $P$ 的直线都与平面 $\alpha$ 平行, 这样的直线有无数条.
        \end{jieda}
        \item 已知两直线 $a//b,a,b \subset$ 平面 $\alpha$ ,直线 $c$ 与 $a$ 异面, $c \cap  b = \varnothing$ ,则 $c$ 与 $\alpha$ 的位置关系是\blkda[2]{平行或相交}
        \begin{jieda}
            直线和平面的关系有 3 种:平行、相交或直线在平面上. 由于直线 $c$ 与 $a$ 异面,故直线不可能在平面上. 在正方体 ${ABCD} - {A}^{\prime }{B}^{\prime }{C}^{\prime }{D}^{\prime }$ 中,取 $a = {AB},b = {CD},\alpha  = {ABCD}$ ,若 $c = {B}^{\prime }{C}^{\prime }$ ,则 $c//\alpha$ ; 取 ${A}^{\prime }{D}^{\prime }$ 和 ${BC}$ 中点 $F,G$ , 若 $c = {FG}$ ,则 $c$ 与 $\alpha$ 相交.
        \end{jieda}
    \end{enumerate}
    \item 在以下四图中,直线 $a$ 与直线 $b$ 可能平行的位置关系只能是\dotfill{(\kh{C})}
    \xx{\includegraphics[width=0.28\textwidth]{bodyxb/src/0_335_275_335_150_0.jpg}}{\includegraphics[width=0.28\textwidth]{bodyxb/src/0_868_276_335_148_0.jpg}}{\includegraphics[width=0.28\textwidth]{bodyxb/src/0_334_450_332_146_0.jpg}}{\includegraphics[width=0.28\textwidth]{bodyxb/src/0_870_449_334_148_0.jpg}}
    \begin{jieda}
        选项 $A$ 中,平面 $\alpha ,\beta$ 内的两直线异面,则 $a$ 与 $b$ 异面;

        选项 $B$ 中,平面 $\alpha ,\beta$ 内的两直线异面,则 $a$ 与 $b$ 异面;

        选项 ${C}$ 中,平面 $\alpha ,\beta$ 内的两直线相交,两相交直线能确定一个平面,则 $a$ 与 $b$ 有可能平行;

        选项 $D$ 中,平面 $\alpha ,\beta$ 内的两直线异面,则 $a$ 与 $b$ 异面.
    \end{jieda}
    \item 设 $\alpha ,\beta$ 是两个不同的平面, $m,n$ 是两条不同的直线,则下列命题正确的是\dotfill{(\kh{B})}
    \xx{若 $m \subset  \alpha ,n \subset  \beta ,\alpha //\beta$ ,则 $m//n$}{若 $m//\alpha ,m \subset  \beta ,\alpha  \cap  \beta  = n$ ，则 $m//n$}{若 $m \subset  \alpha ,n \subset  \alpha ,m//\beta ,n//\beta$ ,则 $\alpha //\beta$}{若 $m \subset  \alpha ,m//\beta ,n \subset  \beta ,n//\alpha$ ,则 $\alpha //\beta$}
    \begin{jieda}
        若 $m \subset  \alpha ,n \subset  \beta ,\alpha //\beta$ ,则 $m//n$ 或异面,故 ${A}$ 错误;

        若 $m//\alpha ,m \subset  \beta ,\alpha  \cap  \beta  = n$ ,由线面平行性质定理可知 $m//n$ ,故 ${B}$ 正确;

        若 $m \subset  \alpha ,n \subset  \alpha ,m//\beta ,n//\beta$ ,当 $m//n$ 时, $\alpha ,\beta$ 可以相交,故 ${C}$ 错误;

        若 $m \subset  \alpha ,m//\beta ,n \subset  \beta ,n//\alpha$ ,当 $m//n$ 时, $\alpha ,\beta$ 可以相交,故 ${D}$ 错误.
    \end{jieda}
\end{enumerate}
\subsection*{证明}
\begin{enumerate}
    \item 在《九章算术》中，四个面都是直角三角形的三棱锥被称为鳖臑. 如图，在鳖臑 $P - {ABC}$ 中, $D$ 为 ${PC}$ 的中点, $M,N$ 分别是 ${AD},{AB}$ 的中点. 证明: ${MN}//$ 平面 ${PBC}$
    \begin{center}
    \includegraphics[width=0.2\textwidth]{bodyxb/src/2_248_212_316_314_0.jpg}
    \end{center}
    \begin{jieda}
        因为 $M,N$ 分别是 ${AD},{AB}$ 的中点,所以 ${MN}//{BD}$，又 ${BD} \subset$ 平面 ${PBC},{MN} \not \subset$ 平面 ${PBC}$ ,所以 ${MN}//$ 平面 ${PBC}$.
    \end{jieda}

    \item 如图,四棱锥 $P - {ABCD}$ 中, ${ABCD}$ 是平行四边形, $M$ 是 ${PC}$ 的中点.
    \begin{enumerate}
        \item 若 ${AB}$ 的中点为 $N$ ，求证: ${MN}//$ 平面 ${APD}$
        \item 在 ${DM}$ 上取一点 $G$ ，过 $G$ 和 ${AP}$ 作平面交平面 ${BDM}$ 于 ${GH}$ ， $H$ 在 ${BD}$ 上，证明: ${AP}//{GH}$
    \end{enumerate}
    \begin{center}
    \includegraphics[width=0.2\textwidth]{bodyxb/src/0_1168_71_368_340_0.jpg}
    \end{center}
    \begin{jieda}
        \begin{enumerate}
            \item 证明如下:

            如图所示,取 ${PD}$ 的中点 $E$ ,连接 ${ME},{AE}$ ,

            因为 $M$ 为 ${PC}$ 的中点,可得 ${ME}//{CD}$ 且 ${ME} = \dfrac{1}{2}{CD}$ ,

            又因为 ${ABCD}$ 为平行四边形,可得 ${AB}//{CD}$ 且 ${AB} = {CD}$ ,

            所以 ${ME}//{AB}$ 且 ${ME} = \dfrac{1}{2}{AB}$ ,

            又因为 $N$ 为 ${AB}$ 的中点,可得 ${ME}//{AN}$ 且 ${ME} = {AN}$ ,

            所以四边形 ${ANME}$ 为平行四边形,所以 ${MN}//{AE}$ ,

            因为 ${MN} \not \subset$ 平面 ${PAD}$ ,且 ${AE} \subset$ 平面 ${PAD}$ ,所以 ${MN}//$ 平面 ${PAD}$ .
            \item 证明如下:

            连接 ${AC}$ 与 ${BD}$ 交于点 $O$ ,且 $O$ 为 ${AC}$ 的中点,

            由点 $M$ 为 ${PC}$ 的中点,所以 ${AP}//{MO}$ ,

            因为 ${AP} \not \subset$ 平面 ${BDM}$ ,且 ${MO} \subset$ 平面 ${BDM}$ ,所以 ${AP}//$ 平面 ${BDM}$ ,

            又因为 ${AP} \subset$ 平面 ${AHGP}$ ,且平面 ${AHGP} \cap$ 平面 ${BDM} = {GH}$ ,所以 ${AP}//{GH}$ .

            \begin{center}
            \includegraphics[width=0.2\textwidth]{bodyxb/src/0_91_1303_366_338_0.jpg}
            \end{center}
        \end{enumerate}
    \end{jieda}

    \item 已知两个全等的正方形 ${ABCD}$ 和 ${ABEF}$ 所在平面相交于 ${AB}$,点 $M$ 、 $N$ 分别在线段 ${AC}$ 、 ${FB}$ 上(与端点不重合),且 ${AM} = {FN}$. 求证: ${MN}//$ 平面 ${BCE}$.

    \begin{center}
        \includegraphics[width=0.2\linewidth]{bodyxb//src/asdfhasldf.png}
    \end{center}

    \begin{jieda}
    \begin{center}
        \includegraphics[width=0.2\linewidth]{bodyxb//src/fa;skjdhfa'hfa.png}
    \end{center}   
    作 ${MP}//{AB}$ 交 ${BC}$ 于点 $P,{NQ}//$  ${AB}$ 交 ${BE}$ 于点 $Q$,所以 ${MP}//{NQ}$.
    
    因为 ${AM} = {FN}$,所以 ${MP} = \dfrac{\sqrt{2}}{2}{MC} = \dfrac{\sqrt{2}}{2}{BN} = {NQ}$. 所以 ${MP} // {NQ}$且$MP = NQ$. 所以四边形 ${MNQP}$ 为平行四边形,所以 ${MN}//{PQ}$.
    
    因为 ${MN} \not \subset$ 平面 ${BCE},{PQ} \subset$ 平面 ${BCE}$,所以 ${MN}//$ 平面 ${BCE}$.
    \end{jieda}

    \item 已知 $P,Q$ 是单位正方体 ${ABCD} - {A}_{1}{B}_{1}{C}_{1}{D}_{1}$ 的面 $A{A}_{1}{D}_{1}D$ ,面 ${A}_{1}{B}_{1}{C}_{1}{D}_{1}$ 的中心 (如图).
    \begin{center}
    \includegraphics[width=0.2\textwidth]{bodyxb/src/bo_d20aucref24c73anc48g_140_392_1501_253_253_0.jpg}
    \end{center}
    \begin{enumerate}
        \item 求线段 ${PQ}$ 的长
        \item 求证: ${PQ}//$ 平面 $A{A}_{1}{B}_{1}B$ 
    \end{enumerate}
    \begin{jieda}
        \begin{enumerate}
            \item 连接 ${A}_{1}{C}_{1}$ ,与 ${B}_{1}{D}_{1}$ 的交点即点 $Q$ . 在 $\bigtriangleup D{A}_{1}{C}_{1}$ 中, ${PQ}$ 是 $D{C}_{1}$ 边上的中位线,故 ${PQ}//D{C}_{1}$ 且 ${PQ} =$  $\dfrac{1}{2}D{C}_{1} = \dfrac{\sqrt{2}}{2}$ . 
            \item 从上一问可知${PQ}//D{C}_{1}$ ,而 $D{C}_{1}//A{B}_{1}$ ,故 ${PQ}//A{B}_{1}$ ,进而 ${PQ}//$ 平面 $A{A}_{1}{B}_{1}B$ ,证毕.
        \end{enumerate}
    \end{jieda}
    \item 用平行于空间四边形 $ {{ABCD}} $ 的一组对边 $ {{AB}}\text{、}{{CD}} $ 的平面截此空间四边形。
    \begin{enumerate}
        \item 求证：所得截面 $ {{MNPQ}} $ 是平行四边形；
        \item 如果 $ {{AB}} = {{CD}} = a $ ,求证：四边形 $ {{MNPQ}} $ 的周长为定值.
    \end{enumerate}
  
    \begin{center}
        \includegraphics[width=0.25\linewidth]{bodyxb//src/cjshf'kashdfws.png}
    \end{center}
    
    \begin{jieda}
        \begin{enumerate}
            \item AB//平面MNPQ，$AB\subset 平面ABC, 平面ABC\cap 平面MNPQ=MN$，由线面平行的判定定理可知，$AB//MN$，同理可得$AB//PQ, CD//MQ, CD//NP$，因此$PQ//MN,NP//MQ$，因此为平行四边形.
            \item 因为$AB//MN,MQ//CD$，因此$\dfrac{BN}{NC}=\dfrac{AM}{MC}=\dfrac{AQ}{QD}=\dfrac{1-t}{t}$，从而周长$=2(ta+(1-t)a)=2a$.
        \end{enumerate}
    \end{jieda}
\end{enumerate}
\subsection*{由线面平行求线段长度}
\begin{enumerate}
    \item 如图，在正方体 ${ABCD} - {A}_{1}{B}_{1}{C}_{1}{D}_{1}$ 中， ${AB} = 2$ ， $E$ 为 ${AD}$ 的中点，点 $F$ 在 ${CD}$ 上，若 ${EF}//$ 平面 $A{B}_{1}C$ ，则 ${EF} =$ \blkda{$\sqrt{2}$}.

    \begin{center}
    \includegraphics[width=0.2\textwidth]{bodyxb/src/0_89_1532_356_348_0.jpg}
    \end{center}
 
    \begin{jieda}
        根据题意,因为 ${EF}//$ 平面 $A{B}_{1}C,{EF} \subset$ 平面 ${ABCD}$ ,且平面 ${ABCD} \cap$ 平面 $A{B}_{1}C = {AC}$ 所以 ${EF}//{AC}$ .

        又 $E$ 是 ${AD}$ 的中点,所以 $F$ 是 ${CD}$ 的中点. 因为在 $\mathrm{{Rt}}\bigtriangleup {DEF}$ 中, ${DE} = {DF} = 1$ ,故 ${EF} = \sqrt{2}$.
    \end{jieda}

    \item 如图,长方体 ${ABCD} - {A}_{1}{B}_{1}{C}_{1}{D}_{1}$ 中, $D{D}_{1} = 8,E,F$ 分别是侧棱 $A{A}_{1},C{C}_{1}$ 上的动点, ${AE} + {CF} = 8$ ，点 $P$ 在棱 $A{A}_{1}$ 上，且 ${AP} = 2$ ，若 ${EF}//$ 平面 ${PBD}$ ，则 ${CF} =$ \blkda{2}.

    \begin{center}
    \includegraphics[width=0.2\textwidth]{bodyxb/src//0_240_626_337_399_0.jpg}
    \end{center}
    \begin{jieda}
        连接 ${AC}$ ,交 ${BD}$ 于点 $O$ ,连接 ${PO}$ .

        \begin{center}
        \includegraphics[width=0.2\textwidth]{bodyxb/src//0_243_1366_336_393_0.jpg}
        \end{center}

        因为 ${EF}//$ 平面 ${PBD},{EF} \subset$ 平面 ${EACF}$ ,平面 ${EACF} \cap$ 平面 ${PBD} = {PO}$ ,所以 ${EF}//{PO}$

        在 $P{A}_{1}$ 上截取 ${PQ} = {AP} = 2$ ,连接 ${QC}$ ,则 ${QC}//{PO}$ ,所以 ${EF}//{QC}$，所以易知四边形 ${EFCQ}$ 为平行四边形,则 ${CF} = {EQ}$.

        又 ${AE} + {CF} = 8,{AE} + {A}_{1}E = 8$ ,所以 ${A}_{1}E = {CF} = {EQ} = \dfrac{1}{2}{A}_{1}Q = 2$ ,故 ${CF} = 2$.
    \end{jieda}
    
    \item 如图，在 $\bigtriangleup  {ABC}$ 中， ${BC} = 9,{BC}//$ 平面 $\alpha$ ，且面 ${ABC} \cap  \alpha  = {MN}$ ，若 $\bigtriangleup  {ABC}$ 的重心 $G$ 在 ${MN}$ 上，则 ${MN} =$ \blkda{6}

    \begin{center}
    \includegraphics[width=0.2\textwidth]{bodyxb/src//1_241_407_345_319_0.jpg}
    \end{center}
    \begin{jieda}
        $\because {BC}//\alpha$ ,且面 ${ABC} \cap  \alpha  = {MN},{BC} \subset$ 面 ${ABC}$ ,

        $\therefore {BC}//{MN}$ ，又重心 $G \in  {MN}$ ，由重心的性质知: $\dfrac{AM}{AB} = \dfrac{MN}{BC} = \dfrac{2}{3}$ ，

        $\therefore {MN} = \dfrac{2}{3}{BC} = 6$.
    \end{jieda}



    \item 如图所示，在棱长为 1 的正方体 ${ABCD} - {A}_{1}{B}_{1}{C}_{1}{D}_{1}$ 中，设 $M$ ， $N$ 分别是线段 $D{A}_{1}$ 、 ${B}_{1}{D}_{1}$ 上的动点，若 ${MN}//$ 平面 $C{C}_{1}{D}_{1}D$ ，则线段 ${MN}$ 长的最小值为\blkda{$\dfrac{\sqrt{2}}{2}$}.

    \begin{center}
    \includegraphics[width=0.2\textwidth]{bodyxb/src/0_90_315_359_357_0.jpg}
    \end{center}
 
    \begin{jieda}
        过点 $M,N$ 分别作 ${MG}//{A}_{1}{D}_{1}$ 交 $D{D}_{1}$ 于点 $G,{NP}//{A}_{1}{D}_{1}$ 交 ${C}_{1}{D}_{1}$ 于点 $P$ ，连接 ${PG}$

        要想 ${MN}//$ 平面 $C{C}_{1}{D}_{1}D$ ,则四边形 ${MGPN}$ 为平行四边形,故 ${NP} = {MG}$ ,

        \begin{center}
        \includegraphics[width=0.2\textwidth]{bodyxb/src/0_1064_924_362_357_0.jpg}
        \end{center}
        

        设 ${D}_{1}G = m \in  \left( {0,1}\right)$ ,则 $P{C}_{1} = m$ ,故 $P{D}_{1} = 1 - m$ ,

        由勾股定理得 ${MN} = {PG} = \sqrt{{D}_{1}{G}^{2} + {D}_{1}{P}^{2}} = \sqrt{{m}^{2} + {\left( 1 - m\right) }^{2}}$ ,

        其中 ${m}^{2} + {\left( 1 - m\right) }^{2} = 2{m}^{2} - {2m} + 1 = 2{\left( m - \dfrac{1}{2}\right) }^{2} + \dfrac{1}{2} \geq  \dfrac{1}{2}$ ,

        当且仅当 $m = \dfrac{1}{2}$ 时,等号成立，故 ${MN} \geq  \dfrac{\sqrt{2}}{2}$.
    \end{jieda}
\end{enumerate}
\subsection*{根据线面平行判断线段比例或点所在位置}
\begin{enumerate}
    \item 如图,已知点 $P$ 在平行四边形 ${ABCD}$ 所在平面外, $E$ 为线段 ${AD}$ 上靠近 $A$ 的三等分点, $F$ 为线段 ${PC}$ 上一点,当 ${PA}//$ 平面 ${EBF}$ 时, $\dfrac{PC}{PF} =$ \blkda{4}
    \begin{center}
    \includegraphics[width=0.2\textwidth]{bodyxb/src/bo_d4lnl5ref24c739do320_3_241_1515_352_258_0.jpg}
    \end{center}
    \begin{jieda}
        如图,连结 ${AC}$ ,交 ${BE}$ 于点 $M$ ,连结 ${FM}$
        
        \begin{center}
        \includegraphics[width=0.2\textwidth]{bodyxb/src/bo_d4lnl5ref24c739do320_4_240_477_352_256_0.jpg}
        \end{center}

        因为 ${PA}//$ 平面 ${EBF}$ ,且 ${PA} \subset$ 平面 ${PAC}$ ,平面 ${PAC} \cap$ 平面 ${BEF} = {FM}$，所以$PA//FM$
        
        因为 ${AE}//{BC}$ ，且 ${AE} : {BC} = 1 : 3$ ，所以 ${AM} : {MC} = 1 : 3$ ，即 ${AM} : {AC} = 1 : 4$，所以 $\dfrac{PC}{PF} = 4$ .
    \end{jieda}


    \item 如图,已知三棱柱 ${ABC} - {A}_{1}{B}_{1}{C}_{1}$ 中， $E$ 是 ${BC}$ 上的动点， $D$ 是 $A{A}_{1}$ 上的动点，且 $\dfrac{AD}{D{A}_{1}} = m$ ， ${AE}//$ 平面 $D{B}_{1}C$ .

    \begin{center}
    \includegraphics[width=0.2\textwidth]{bodyxb/src/1_90_1754_300_348_0.jpg}
    \end{center}

    \begin{enumerate}
        \item 若 $E$ 是 ${BC}$ 的中点，则 $m$ 的值为\blkda{1}
        \item 若 $E$ 是 ${BC}$ 上靠近 $B$ 的三等分点，则 $m$ 的值为\blkda{2}
    \end{enumerate}
    \begin{jieda}
        \begin{enumerate}
            \item 如图,设平面 ${ADE}$ 与直线 ${B}_{1}C$ 相交于点 $G$ ,连接 ${DG},{GE}$ .
            
            因为 ${AE}//$ 平面 $D{B}_{1}C,{AE} \subset$ 平面 ${ADE}$ ，平面 ${ADE} \cap$ 平面 $D{B}_{1}C = {DG}$ ， 所以 ${AE}//{DG}$.

            又 ${AD}//B{B}_{1}$ ,所以 ${AD}//$ 平面 ${CB}{B}_{1}{C}_{1}$ ,

            又 ${AD} \subset$ 平面 ${ADE}$ ，平面 ${ADE} \cap$ 平面 ${CB}{B}_{1}{C}_{1} = {EG}$，

            所以 ${AD}\parallel {EG}$ ,则四边形 ${DAEG}$ 是平行四边形.

            故 ${DA} = {GE} = \dfrac{1}{2}{B{B}_{1}}$ ，所以 $G$ 是 $C{B}_{1}$ 的中点.

            故 ${AD} = D{A}_{1}$ ，即 $\dfrac{AD}{D{A}_{1}} = 1$ ，即 $m = 1$.
            \begin{center}
            \includegraphics[width=0.2\textwidth]{bodyxb/src/2_1176_296_299_353_0.jpg}
            \end{center}
            \item 如图，设平面 ${ADE}$ 与直线 ${B}_{1}C$ 相交于点 $H$ ，连接 ${DH}$ ， ${HE}$ .

            因为 ${AE}//$ 平面 $D{B}_{1}C$ ， ${AE} \subset$ 平面 ${ADE}$ ，平面 ${ADE} \cap$ 平面 $D{B}_{1}C = {DH}$ ，

            \begin{center}
            \includegraphics[width=0.2\textwidth]{bodyxb/src/2_1203_971_291_337_0.jpg}
            \end{center}
            

            所以 ${AE}//{DH}$ ，又 ${AD}//{B{B}_{1}}$ ，所以 ${AD}//$ 平面 ${CB}{B}_{1}{C}_{1}$，

            又 ${AD} \subset$ 平面 ${ADE}$ ，平面 ${ADE} \cap$ 平面 ${CB}{B}_{1}{C}_{1} = {EH}$，

            所以 ${AD}//{EH}$ ,故四边形 ${DAEH}$ 是平行四边形,

            则 ${AD} = {EH}$ ,且 ${EH}//B{B}_{1}$,

            所以 $\dfrac{EH}{B{B}_{1}} = \dfrac{CE}{CB} = \dfrac{2}{3}$ ,所以 $\dfrac{AD}{A{A}_{1}} = \dfrac{EH}{B{B}_{1}} = \dfrac{2}{3}$，则 $\dfrac{AD}{D{A}_{1}} = 2$ ,即 $m = 2$.
        \end{enumerate}
    \end{jieda}
    \item 如图,在三棱柱 ${ABC} - {A}_{1}{B}_{1}{C}_{1}$ 中, $E$ 是棱 $C{C}_{1}$ 上的一点,且 $\dfrac{{C}_{1}E}{EC} = \dfrac{2}{3},D$ 是棱 ${BC}$ 上一点. 若 ${A}_{1}B//$ 平面 ${ADE}$ ,则 $\dfrac{BD}{BC}$ 的值为\blkda{$\dfrac{5}{8}$}.

    \begin{center}
    \includegraphics[width=0.2\textwidth]{bodyxb/src/3_91_1525_312_371_0.jpg}
    \end{center}

    \begin{jieda}
        连接 ${A}_{1}C\text{ 、 }{AE}$ 相交于点 $O$ ,连接 ${OD}$ ,

        因为 ${A}_{1}B//$ 平面 ${ADE}$ ,平面 ${A}_{1}{BC} \cap$ 平面 ${ADE} = {OD},{A}_{1}B \subset$ 平面 ${A}_{1}{BC}$ , 所以 ${A}_{1}B//{OD}$ ,所以 $\dfrac{CO}{{A}_{1}O} = \dfrac{DC}{BD}$ ,

        因为 $A{A}_{1}//{CE}$ ,所以 $\bigtriangleup {A}_{1}O{A}^{CD}\bigtriangleup {COE}$ ,

        所以 $\dfrac{CO}{{A}_{1}O} = \dfrac{CE}{A{A}_{1}} = \dfrac{3}{5}$ ,即 $\dfrac{DC}{BD} = \dfrac{3}{5}$，可得 $\dfrac{BD}{BC} = \dfrac{5}{8}$.
        \begin{center}
        \includegraphics[width=0.2\textwidth]{bodyxb/src/3_1124_1832_315_372_0.jpg}
        \end{center}
    \end{jieda}
    \item 四棱锥 $P - {ABCD}$ 的底面为平行四边形，如图所示，点 $E$ 是棱 ${PD}$ 上一点， ${PE} = \dfrac{2}{3}{PD}$ ，若 $\vv{PF} = \lambda \vv{PC}$ 且满足 ${BF}//$ 平面 ${ACE}$ ，则 $\lambda  =$ \blkda{$\dfrac{1}{2}$}.

    \begin{center}
    \includegraphics[width=0.2\textwidth]{bodyxb/src/4_90_477_410_377_0.jpg}
    \end{center}
    \begin{jieda}
        如图,连接 ${BD}$ ,交 ${AC}$ 于点 $O$ ,连接 ${OE}$ ,

        由四边形 ${ABCD}$ 是平行四边形,得 ${BO} = {OD}$ ,

        在线段 ${PE}$ 上取点 $G$ ,使得 ${GE} = {ED}$ ,由 ${PE} = \dfrac{2}{3}{PD}$ ，得 $\dfrac{PG}{PE} = \dfrac{1}{2}$ ，

        连接 ${BG},{FG}$ ,则 ${BG}//{OE}$ ,由 ${OE} \subset$ 平面 ${ACE},{BG} \text{ ⊄ }$ 平面 ${ACE}$ ,

        得 ${BG}//$ 平面 ${ACE}$ ,而 ${BF}//$ 平面 ${ACE},{BG} \cap  {BF} = B,{BG},{BF} \subset$ 平面 ${BGF}$ ,

        因此平面 ${BGF}//$ 平面 ${ACE}$ ,又平面 ${PCD} \cap$ 平面 ${ACE} = {EC}$ ,平面 ${PCD} \cap$ 平面 ${BGF} = {GF}$ ,

        则 ${GF}//{EC}$ ,所以 $\lambda  = \dfrac{PF}{PC} = \dfrac{PG}{PE} = \dfrac{1}{2}$ .

        \begin{center}
        \includegraphics[width=0.2\textwidth]{bodyxb/src/4_88_1363_429_385_0.jpg}
        \end{center}
    \end{jieda}
\end{enumerate}

\section*{线面垂直}
\subsection*{位置关系}
\begin{enumerate}
    \item 若直线 $a \bot$ 直线 $b$ ,且 $a \bot$ 平面 $\alpha$ ,则有\dotfill{(\kh{D})}
    \xx{$b//\alpha$}{$b \subset  \alpha$}{$b \bot  \alpha$}{$b//\alpha$ 或 $b \subset  \alpha$}
    \begin{jieda}
        设直线 $a$ 与平面 $\alpha$ 垂直相交于点 $O$ ,并有引理: 过点 $O$ 与直线 $a$ 垂直的直线 $c$ 必在平面 $\alpha$ 上 (证明见注释). 当 $b \subset  \alpha$ 时, $a \bot  b$ 自然成立. 当 $b \subset  \alpha$ 时,直线 $b$ 必不过点 $O$ ,此时直线 $b$ 和点 $O$ 确定一个平面 $\beta$ ,在平面 $\beta$ 内,过点 $O$ 作直线 ${b}^{\prime }//b$ ,则 $b\bot a \Longrightarrow  {b}^{\prime }\bot a$ ,根据引理, ${b}^{\prime }\square \alpha$ ,再由直线与平面平行的判定定理,得 $b//a$ . 注: 相交直线 $a,c$ 确定一个平面,记为 $\gamma$ ；根据公理3,平面 $\alpha ,\gamma$ 的交线 ${c}^{\prime }$ 过点 $O$ ；此时,在平面 $\gamma$ 内,直线 $c$ 和 ${c}^{\prime }$ 均与直线 $a$ 垂直,故 $c = {c}^{\prime }$ ,从而 $c \subset  \alpha$ ,证毕.
    \end{jieda}
    \item 若直线 $a//$ 平面 $\alpha$ ,直线 $b\bot$ 平面 $\alpha$ ,则 $a$ 与 $b$ 不可能\dotfill{(\kh{C})}
    \xx{相交}{异面}{平行}{垂直}
    \begin{jieda}
        事实上, $a,b$ 一定互相垂直: 过 $a$ 作平面. 与 $a$ 交于 ${a}^{\prime }$ ,则 $b\bot {a}^{\prime } \Longrightarrow  b\bot a$ . 当 $a,b$ 共面时,两者相交; 当 $a,b$ 不共面时,根据定义,两者异面; 而 $a//b$ 必不可能.
    \end{jieda}
    
    \item 异面直线在同一平面上的射影不可能是\dotfill{(\kh{B})}
    \xx{两平行直线}{同一直线}{两相交直线}{一点与一直线}
    \begin{jieda}
        在正方体 ${ABCD} - {A}^{\prime }{B}^{\prime }{C}^{\prime }{D}^{\prime }$ 中. 取平面 $- {ABCD}$ ,两异面直线取 ${A}^{\prime }{D}^{\prime }$ 和 $B{C}^{\prime }$ 时,即得选项 ${A}_{1}$ 取 ${A}^{\prime }{D}^{\prime }$ 和 $A{C}^{\prime }$ 时，即得选项 $C$ ; 取 ${A}^{\prime }{D}^{\prime }$ 和 $B{B}^{\prime }$ 时,即得选项 $D$ . 选项 $B$ 是不可能的,证明如下: 假设直线 ${l}_{1},{l}_{2}$ 在平面 $\alpha$ 上的射影为同一直线,在该射影上任取两点 $P,Q$ ,过这两点作平面 $\alpha$ 的垂线 ${m}_{1},{m}_{2}$ ,根据射影的定义, ${m}_{1},{m}_{2}$ 与 ${l}_{1},{l}_{2}$ 均有交点. 又 ${m}_{1}//$  ${m}_{2}$ ,故 ${m}_{1},{m}_{2}$ 确定一个平面,由公理 $1,{l}_{1},{l}_{2}$ 也在这个平面内,故 ${l}_{1},{l}_{2}$ 共面,与题干矛盾.
    \end{jieda}
    
    \item 已知 ${PH}$ 垂直于 ${{Rt}} \bigtriangleup  {HEF}$ 所在的平面,且 ${HE} \bot  {EF}$ ,连接 ${PE},{PF}$ , 则图中直角三角形的个数是\dotfill{(\kh{D})}
    
    \begin{center}
    \includegraphics[width=0.2\textwidth]{bodyxb/src/bo_d20aucref24c73anc48g_142_1267_1894_242_197_0.jpg}
    \end{center}

    \xx{1}{2}{3}{4}
    \begin{jieda}
        不难判断, $\bigtriangleup {HEF},\bigtriangleup {PHE},\bigtriangleup {PHF}$ 均为直角三角形. 又 ${PH} \bot$ 平面 ${HEF} \Longrightarrow  {PH} \bot  {EF}$ ,而 ${EF} \bot  {HE} \cap$  ${PH} = H,\therefore {EF} \bot$ 平面 ${PHE},\therefore {EF} \bot  {PE}$ ,即 $\bigtriangleup  {PEF}$ 也是直角三角形.
    \end{jieda}
    
    \item 若两直线 $a$ 与 $b$ 异面,则过 $a$ 且与 $b$ 垂直的平面 \dotfill{(\kh{B})}
    \xx{有且只有一个}{可能存在也可能不存在}{有无数多个}{一定不存在}
    \begin{jieda}
        假设存在这样一个平面 $\alpha$ 满足要求, $\because b\bot \alpha ,a \subset  \alpha ,\therefore b\bot a$ . 因此 $a$ 与 $b$ 互相垂直是平面 $\alpha$ 存在的必要条件. 而当 $a\bot b$ 时,由 $a,b$ 的公垂线和 $a$ 确定的平面,即为 $\alpha$ . 
        
        当 $a\bot b$ 时,我们已说明了平面 $\alpha$ 的存在性,以下说明其唯一性. 假设平面 $\alpha ,\beta$ 均满足条件, $\alpha ,\beta$ 与直线 $b$ 交于相异两点 $A,B$ . 在直线 $a$ 上任取一点 $P$ ,则在平面 ${PAB}$ 中, ${PA}\bot b$ , ${PB} \bot  b$ ,矛盾. 故平面 $\alpha$ 是唯一的.
    \end{jieda}
    
    \item 从平面 $\alpha$ 外一点 $P$ 引与 $\alpha$ 相交的直线,使得点 $P$ 与交点的距离为 1,则满足条件的直线一定不可能是\dotfill{(\kh{C})}
    \xx{0 条}{1 条}{2 条}{无数条}
    \begin{jieda}
        记点 $P$ 与平面 $\alpha$ 的距离为 $d$ . 若 $d > 1$ ,则满足条件的直线不存在; 若 $d = 1$ ,则仅有一条; 若 $d < 1$ ,则有无数条.
    \end{jieda}
    
    \item 已知 $a,b$ 为异面直线,则下列命题正确的是\dotfill{(\kh{C})}
    \xx{过 $a,b$ 外一点 $P$ 一定可以引一条与 $a,b$ 都平行的直线}{过 $a,b$ 外一点 $P$ 一定可以作一个与 $a,b$ 都平行的平面}{过 $a$ 一定可以作一个与 $b$ 平行的平面}{过 $a$ 一定可以作一个与 $b$ 垂直的平面}

    \begin{jieda}
        对于选项 A, 若这样的直线存在, 由公理 4 可知: $a // b$, 与题设相悖, 错误; 对于选项 B, 当点P在过$a$且与$b$平行的直线上则无法实现, 错误; 对于选项 C,存在性: 在直线 $a$ 上任取一点 $P$ ,在直线 $b$ 上任取两点 $Q,R$ ,在平面 ${PQR}$ 中 (即点 $P$ 和直线 $b$ 确定的平面), 过点 $P$ 作直线 $b$ 的平行线 ${b}^{\prime }$ (可由尺规作图,构造平行四边形得到). 此时,相交直线 ${b}^{\prime }$ 和 $a$ 确定的平面即为与 $b$ 平行的平面 $\alpha$，正确; 对于选项 $D$，若$a$与$b$不垂直的时候无法实现，错误.
    \end{jieda}
\end{enumerate}

\subsection*{证明}
\begin{enumerate}
    \item \begin{enumerate}
        \item 已知 ${AB},{CD}$ 是两条不在同一个平面上的线段,且 ${AC} = {AD},{BC} = {BD}$ (如图),求证: ${AB}\bot {CD}$
        \begin{center}
        \includegraphics[width=0.2\textwidth]{bodyxb/src/bo_d20aucref24c73anc48g_143_417_1279_278_254_0.jpg}
        \end{center}
        \begin{jieda}
            取$CD$中点$O$，连接$AO, BO$，因为$AC = AD$，则$\triangle ACD$为等腰三角形，所以$AO \bot CD$，同理$BO \bot CD$，因此$CD \bot$ 平面$AOB$，故$AB \bot CD$.
        \end{jieda}
        \item 在空间四边形 ${ABCD}$ 中, $E$ , $F$ 分别为 ${AD}$ , ${BC}$ 的中点,若 ${AC} = {BD} = a,{EF} = \dfrac{\sqrt{2}}{2}a$ , $\angle {BDC} = {90}^{ \circ  }$ (如图),求证: ${BD} \bot$ 平面 ${ACD}$
        \begin{center}
        \includegraphics[width=0.2\textwidth]{bodyxb/src/bo_d20aucref24c73anc48g_143_953_1284_258_245_0.jpg}
        \end{center}
        \begin{jieda}
            取$CD$边中点$O$，连接$EO, FO$，根据三角形中位线定理可知$EO = FO = \dfrac{a}{2}$，故根据勾股定理逆定理可知$\angle{EOF}$为直角，即$EO \bot FO$，因此$EO \bot BD$，又因为$BD \bot CD$，$CD \cap EO = O$，所以${BD} \bot$ 平面 ${ACD}$
        \end{jieda}
        \item 在正方体 ${ABCD} - {A}_{1}{B}_{1}{C}_{1}{D}_{1}$ 中, $O$ 为底面 ${ABCD}$ 的中心, ${B}_{1}H \bot  {D}_{1}O$ , $H$ 是垂足 (如图),求证: ${B}_{1}H \bot$ 平面 $A{D}_{1}C$
        
        \begin{center}
        \includegraphics[width=0.2\textwidth]{bodyxb/src/bo_d20aucref24c73anc48g_143_435_1845_255_256_0.jpg}
        \end{center}
        \begin{jieda}
            因为$AC \bot BD, AC \bot BB_1$，故$AC \bot $平面$BB_1D_1D$，又因为$B_1H \subset $平面$BB_1D_1D$，因此$AC \bot B_1H$，又因为$B_1H \bot D_1O, D_1O \cap AC = O$，因此$B_1H \bot$ 平面 $A{D}_{1}C$
        \end{jieda}
        \item 在正方体 ${ABCD} - {A}_{1}{B}_{1}{C}_{1}{D}_{1}$ 中, $G$ 为 $C{C}_{1}$ 的中点, ${AC}$ 交 ${BD}$ 于点 $O$ (如图),求证: ${A}_{1}O \bot$ 平面 ${GBD}$
        \begin{center}
        \includegraphics[width=0.2\textwidth]{bodyxb/src/bo_d20aucref24c73anc48g_143_948_1841_252_259_0.jpg}
        \end{center}
        \begin{jieda}
            $AC \bot BD, AA_1 \bot BD \Longrightarrow BD \bot $平面$ACC_1A_1 \Longrightarrow BD \bot A_1O$

            设立方体边长为$2$，则$A_1O = \sqrt{6}, OG = \sqrt{3}, A_1G = 3$，故$A_1O \bot OG$

            因为$A_1O \cap BD = O$故${A}_{1}O \bot$ 平面 ${GBD}$
        \end{jieda}
    \end{enumerate}

    \item 在 $\bigtriangleup {ABC}$ 中,已知 $\angle {ABC} = {90}^{ \circ  },{SA} \bot$ 平面 ${ABC}$ ,点 $A$ 在 ${SB},{SC}$ 上的射影分别为 $M,N$ (如图). 求证:

    \begin{enumerate*}
        \item ${BC} \bot$ 平面 ${SAB}$
        \item ${AM} \bot$ 平面 ${SBC}$
        \item ${SC} \bot  {MN}$
    \end{enumerate*}
    \begin{center}
    \includegraphics[width=0.2\textwidth]{bodyxb/src/bo_d20aucref24c73anc48g_144_425_483_251_204_0.jpg}
    \end{center}
    \begin{jieda}
        \begin{enumerate}
            \item 因为$\angle {ABC} = {90}^{ \circ  }$，所以$AB \bot BC$，因为${SA} \bot$ 平面 ${ABC}$，所以$SA \bot BC$，因此${BC} \bot$ 平面 ${SAB}$

            \item 由${BC} \bot$ 平面 ${SAB}$可知$BC \bot AM$，又因为$AM \bot SB$，所以${AM} \bot$ 平面 ${SBC}$

            \item 由${AM} \bot$ 平面 ${SBC}$可知$AM \bot SC$，又因为$AN \bot SC$，因此$SC \bot$ 平面$AMN$，因此${SC} \bot  {MN}$
        \end{enumerate}
    \end{jieda}
\end{enumerate}
\subsection*{三垂线定理}
\begin{enumerate}
    \item 如图所示,在三棱锥 $P - {ABC}$ 中, ${PA} \bot$ 平面 ${PBC},{PB} = {PC},M$ 是 ${BC}$ 的中点,求证: ${BC} \bot  {AM}$ .
    
    \begin{center}
    \includegraphics[width=0.2\textwidth]{bodyxb/src/bo_d4ogc4bef24c73bbtdng_1_1225_223_565_440_0.jpg}
    \end{center}
    \begin{jieda}
        因为$PB = PC, M$ 是 ${BC}$ 的中点所以$PM \bot BC$

        因为${PA} \bot$ 平面 ${PBC}$，所以$AM$在直线内的射影为$PM$，根据三垂线定理可知$BC \bot AM$
    \end{jieda}
    \item 如图所示,点 $O$ 是 $\bigtriangleup {ABC}$ 的垂心, ${OP} \bot$ 平面 ${ABC}$ . 求证: ${PA} \bot  {BC}$ .

    \begin{center}
    \includegraphics[width=0.2\textwidth]{bodyxb/src/bo_d4ogc4bef24c73bbtdng_1_482_1083_610_385_0.jpg}
    \end{center}

    \begin{jieda}
        因为${PO} \bot$ 平面 ${ABC}$，所以$PA$在直线内的射影为$AO$，因为$O$为垂心，所以$BC \bot AO$，故根据三垂线定理可知$BC \bot PA$
    \end{jieda}
    \item 如图所示,在空间四边形 ${ABCD}$ 中,已知 ${AB} \bot  {CD},{BC} \bot  {AD}$ . 求证: ${AC} \bot  {BD}$ .
    
    \begin{center}
    \includegraphics[width=0.2\textwidth]{bodyxb/src/bo_d4ogc4bef24c73bbtdng_2_607_317_489_390_0.jpg}
    \end{center}
    \begin{jieda}
        设$A$在面$BCD$内的射影为$O$，则根据三垂线定理可知$BO \bot CD$，$DO \bot BC$，因此$O$为$\triangle BCD$的垂心，故$CO \bot BD$，根据三垂线定理可知$AC \bot BD$
    \end{jieda}

    \item 如图所示,在正方体 ${ABCD} - {A}_{1}{B}_{1}{C}_{1}{D}_{1}$ 中,点 $P,Q,R$ 分别是线段 ${B}_{1}B,{AB},{A}_{1}C$ 上的动点
        
    \begin{center}
    \includegraphics[width=0.2\textwidth]{bodyxb/src/bo_d4ogc4bef24c73bbtdng_2_1348_173_546_529_0.jpg}
    \end{center}

    观察直线 ${CP}$ 与 ${D}_{1}Q,{CP}$ 与 ${D}_{1}R$ ,给出下列结论:
    \begin{dingautolist}{172}
        \item 存在点 $P$ ，对于任意给定的点 $Q$ ，使得 ${CP}\bot {D}_{1}Q$
        \item 存在点 $Q$ ，对于任意给定的点 $P$ ，使得 ${D}_{1}Q \bot  {CP}$
        \item 存在点 $P$ ,对于任意给定的点 $R$ ,使得 ${CP} \bot  {D}_{1}R$
        \item 存在点 $R$ ,对于任意给定的点 $P$ ,使得 ${D}_{1}R \bot  {CP}$
    \end{dingautolist}
    其中正确的是\blkda{\ding{172}}.
\end{enumerate}
\subsection*{点在面内的射影}
\begin{enumerate}
    \item 已知 ${SA}$ 垂直于直角三角形 ${ABC}$ 所在平面, $\angle {ABC} = {90}^{ \circ  },B{B}_{1}\bot {AC}$ 于点 ${B}_{1}$ ,则:
    \begin{enumerate}
        \item ${SA},{SB},{SC}$ 在平面 ${ABC}$ 上的射影依次是\blkda[3]{点 $A\text{ 、 }{AB}\text{ 、 }{AC}$}
        \item ${CA},{CB},{CS}$ 在平面 ${SAB}$ 上的射影依次是\blkda[3]{${BA}$ 、点 $B$ 、 ${BS}$}
        \item ${BA},{BC},{BS}$ 在平面 ${SAC}$ 上的射影依次是\blkda[3]{${B}_{1}A$ 、 ${B}_{1}C$ 、 ${B}_{1}S$}
    \end{enumerate}
    \item 已知 $P$ 为 $\bigtriangleup {ABC}$ 所在平面外一点,且在平面 ${ABC}$ 上的射影为 $O$ .
    \begin{enumerate}
        \item 若 ${PA} = {PB} = {PC}$ ,则 $O$ 为 $\bigtriangleup  {ABC}$ 的\blkda{外}心,又若 $\angle {ACB} = {90}^{ \circ  }$ ,则 $O$ 在\blkda[2.5]{边 ${AB}$ 的中点}
        \begin{jieda}
            ${PA} = {PB} = {PC} \Longrightarrow  {OA} = {OB} = {OC}$ ,故 $O$ 为 $\bigtriangleup  {ABC}$ 的外心；又直角三角形的外心在斜边的中点,故 $\angle {ACB} = {90}^{ \circ  }$ 时, $O$ 在边 ${AB}$ 的中点.
        \end{jieda}
        \item 若 ${PA},{PB},{PC}$ 与平面 ${ABC}$ 所成的角相等,则 $O$ 为 $\bigtriangleup  {ABC}$ 的\blkda{外}心
        \begin{jieda}
            由题意, $\bigtriangleup {PAO}$ , $\bigtriangleup {PBO}$ , $\bigtriangleup {PCO}$ 是全等三角形,故 ${OA} = {OB} = {OC}$ , $O$ 为 $\bigtriangleup {ABC}$ 的外心.
        \end{jieda}
        \item 若 $O$ 在 $\bigtriangleup {ABC}$ 内,且 $P$ 到 $\bigtriangleup {ABC}$ 的三边距离相等,则 $O$ 为 $\bigtriangleup {ABC}$ 的\blkda{内}心
        \begin{jieda}
            过点 $O$ 作 $\bigtriangleup {ABC}$ 三边 ${AB},{BC},{CA}$ 的垂线,垂足分别为 $D,E,F$ . 由三垂线定理,点 $P$ 到 $\bigtriangleup {ABC}$ 三边 ${AB}$ 、 ${BC}\text{ 、 }{CA}$ 的距离就是 ${PD},{PE},{PF}\text{ 、 }{PD} = {PE} = {PF} \Longrightarrow  {OD} = {OE} = {OF}$ ,又 $O$ 在 $\bigtriangleup {ABC}$ 内,故 $O$ 为 $\bigtriangleup {ABC}$ 的内心. (注: 若不限定 $O$ 在 $\bigtriangleup {ABC}$ 内,则 $O$ 为 $\bigtriangleup {ABC}$ 的内心或旁心)
        \end{jieda}
        \item 若四面体 ${PABC}$ 的两组对棱互相垂直,则 $O$ 为 $\bigtriangleup  {ABC}$ 的\blkda{垂}心
        \begin{jieda}
            连接 ${AO}$ ,它是 ${AP}$ 在平面 ${ABC}$ 上的射影. 由三垂线定理, ${PA} \bot  {BC} \Longrightarrow  {AO} \bot  {BC}$ . 同理 ${CO} \bot  {AB}$ 、 ${BO} \bot  {AC}$ , 故 $O$ 为 $\bigtriangleup {ABC}$ 的垂心.
        \end{jieda}
    \end{enumerate}
    \item 已知 $P$ 是四边形 ${ABCD}$ 所在平面外一点,且 $P$ 到这四边形各边的距离相等,那么这个四边形一定是 \dotfill{(\kh{C})}
    \xx{圆内接四边形}{矩形}{圆外切四边形}{平行四边形}
    \begin{jieda}
        记点 $P$ 在平面 ${ABCD}$ 上的射影是点 $O$ ,由题意及勾股定理,点 $O$ 到四边形 ${ABCD}$ 各边的距离 $r$ 也相等。 $\therefore$ 这个四边形是圆外切四边形,该圆以点 $O$ 为圆心, $r$ 为半径.
    \end{jieda}
    \item 已知 $P$ 是 $\bigtriangleup {EFG}$ 所在平面外一点,且 ${PE} = {PG}$ ,则点 $P$ 在面 ${EFG}$ 上的射影一定在 $\bigtriangleup  {EFG}$ 的\dotfill{(\kh{D})}
      \xx{$\angle {FEG}$ 的平分线上}{边 ${EG}$ 的高上}{边 ${EG}$ 的中线上}{边 ${EG}$ 的垂直平分线上}
      \begin{jieda}
          设点 $P$ 在面 ${EFG}$ 上的射影是点 $O$ ,在 ${{Rt}} \bigtriangleup  {POE}$ 和 ${{Rt}} \bigtriangleup  {POG}$ 中,由勾股定理,易得: ${PE} = {PG} \Longrightarrow  {OE} = {OG}$ ,即点 $O$ 在边 ${EG}$ 的垂直平分线上.
      \end{jieda}
      
\end{enumerate}
\subsection*{点到平面的距离}
\begin{enumerate}
    \item 若正 $\bigtriangleup {ABC}$ 的边长为 $\dfrac{4\sqrt{3}}{3}$ ，则到三个顶点的距离都为 1 的平面有 \dotfill{(\kh{C})}

    \xx{1 个} {3 个} {5 个} {7 个}

    \item 在棱长为 $\alpha$ 的正方体 ${ABCD} - {A}_{1}{B}_{1}{C}_{1}{D}_{1}$ 中， $A$ 到直线 ${B}_{1}C$ 的距离为\blkda{$\dfrac{\sqrt{6}a}{2}$}， $A{A}_{1}$ 到平面 $B{B}_{1}{D}_{1}D$ 的距离为\blkda{$\dfrac{\sqrt{2}a}{2}$}， $A$ 到平面 ${A}_{1}{BD}$ 的距离为\blkda{$\dfrac{\sqrt{3}a}{3}$}.
    

    \item \begin{enumerate}
        \item 在长方体 ${ABCD} - {A}_{1}{B}_{1}{C}_{1}{D}_{1}$ 中,棱 ${AB},{AD},A{A}_{1}$ 的长分别为6,8,3.6,则 $A$ 到 ${B}_{1}{D}_{1}$ 的距离为\blkda{6}
        \begin{jieda}
            在平面 $A{B}_{1}{D}_{1}$ 中,作 ${AE}\bot {B}_{1}{D}_{1}$ 于 $E$ ,连接 ${A}_{1}E$ . 显然, ${AE}$ 在平面 ${A}_{1}{B}_{1}{C}_{1}{D}_{1}$ 上的射影是 ${A}_{1}E$ ,由三垂线定理, ${AE} \bot  {B}_{1}{D}_{1} \Longrightarrow  {A}_{1}E \bot  {B}_{1}{D}_{1}$ . 由面积法 (或相似三角形), ${A}_{1}E = \dfrac{{A}_{1}{D}_{1} \cdot  {A}_{1}{B}_{1}}{{B}_{1}{D}_{1}} = {4.8}$ ,故 ${AE} =$  $\sqrt{{A}_{1}{A}^{2} + {A}_{1}{E}^{2}} = 6$ ,即为所求.
        \end{jieda}
        \item 在矩形 ${ABCD}$ 中,已知 ${AB} = b,{AD} = c$ ,线段 ${PA}\bot$ 矩形 ${ABCD}$ 所在平面,且 ${PA} =$  $a$ ,则 $\bigtriangleup  {PBD}$ 的面积是\blkda[3.5]{$\dfrac{1}{2}\sqrt{{a}^{2}{b}^{2} + {b}^{2}{c}^{2} + {c}^{2}{a}^{2}}$}
        \begin{jieda}
            在平面 ${PBD}$ 中,作 ${PE}\bot {BD}$ 于 $E$ ,连接 ${AE}$ . 显然, ${PE}$ 在平面 ${ABCD}$ 上的射影是 ${AE}$ ,由三垂线定理, ${PE} \bot  {BD} \Longrightarrow  {AE} \bot  {BD}$ . 由面积法 (或相似三角形), ${AE} = \dfrac{{AD} \cdot  {AB}}{BD} = \dfrac{bc}{\sqrt{{b}^{2} + {c}^{2}}}$ ,故 ${PE} =$  $\sqrt{A{P}^{2} + A{E}^{2}} = \sqrt{\dfrac{{a}^{2}{b}^{2} + {b}^{2}{c}^{2} + {c}^{2}{a}^{2}}{{b}^{2} + {c}^{2}}}$ ,故 ${S}_{\bigtriangleup {PBD}} = \dfrac{1}{2} \cdot  {BD} \cdot  {PE} = \dfrac{1}{2}\sqrt{{a}^{2}{b}^{2} + {b}^{2}{c}^{2} + {c}^{2}{a}^{2}}$
        \end{jieda}
        \item 在棱长为 $a$ 的正方体 ${ABCD} - {A}_{1}{B}_{1}{C}_{1}{D}_{1}$ 中, ${D}_{1}$ 到 ${B}_{1}C$ 的距离为\blkda{$\dfrac{\sqrt{6}}{2}a$}, $A$ 到 ${A}_{1}C$ 的距离为\blkda{$\dfrac{\sqrt{6}}{3}a$}
        \begin{jieda}
            $\bigtriangleup {D}_{1}{B}_{1}C$ 是边长为 $\sqrt{2}a$ 的等边三角形,故 ${D}_{1}$ 到 ${B}_{1}C$ 的距离为 $\dfrac{\sqrt{3}}{2} \cdot  \sqrt{2}a = \dfrac{\sqrt{6}}{2}a.\operatorname{Rt} \bigtriangleup  {A}_{1}{AC}$ 中, $\angle {A}_{1}{AC} = {90}^{ \circ  }$ ,由面积法 (或相似三角形), $A$ 到 ${A}_{1}C$ 的距离为 $\dfrac{A{A}_{1} \cdot  {AC}}{{A}_{1}C} = \dfrac{a \cdot  \sqrt{2}a}{\sqrt{3}a} = \dfrac{\sqrt{6}}{3}a.$
        \end{jieda}
        \item 在菱形 ${ABCD}$ 中,已知 $\angle {BAD} = {60}^{ \circ  },{AB} = {10}{{cm}},{PA}\bot$ 菱形 ${ABCD}$ 所在平面,且 ${PA} = {5{{cm}}}$ ,则 $P$ 到 ${BD}$ 的距离为\blkda{10} ${{cm}}$ , $P$ 到 ${DC}$ 的距离为\blkda{10} ${{cm}}$
        \begin{jieda}
            连接 ${AC},{BD}$ 交于点 $E$ ,连接 ${PE}.{AE}$ 是 ${PE}$ 在平面 ${ABCD}$ 上的射影,由三垂线定理, ${BD}\bot {AE} \Longrightarrow  {BD}$  $\bot {PE}.\bigtriangleup {ABD}$ 是等边三角形,故 ${AE} = 5\sqrt{3}{{cm}},P$ 到 ${BD}$ 的距离 ${PE} = \sqrt{A{P}^{2} + A{E}^{2}} = {10}{{cm}}$ . 作 ${AF}\bot {CD}$ 延长线于点 $F$ ,连接 ${PF}.{AF}$ 是 ${PF}$ 在平面 ${ABCD}$ 上的射影,由三垂线定理, ${AF} \bot  {CD} \Longrightarrow  {PF} \bot  {CD}$ . 不难求得, ${AF} = {AE} = 5\sqrt{3}$  ${{cm}}$ ,故 $P$ 到 ${DC}$ 的距离 ${PF} = {PE} = {10}{{cm}}$ .
        \end{jieda}
    \end{enumerate}

    \item 从平面 $\alpha$ 外一点 $P$ 向 $\alpha$ 引垂线和若干条斜线,若这些斜线与 $\alpha$ 所成的角都相等,则\dotfill{(\kh{C})}
    \xx{斜足一定是一正多边形的顶点}{垂足是以斜足为顶点的多边形的内心}{垂足是以斜足为顶点的多边形的外心}{垂足是以斜足为顶点的多边形的垂心}
    \begin{jieda}
        由题意,垂线、斜线段及其对应的射影构成的三角形均全等,故垂足到各个斜足的距离均相等,换言之:垂足在任意两斜足连线段的垂直平分线上,故垂足是以斜足为顶点的多边形的外心.
    \end{jieda}
    \item \begin{enumerate}
        \item 已知 ${PA}$ 垂直于 $\bigtriangleup  {ABC}$ 所在平面,且 ${AB} = {AC} = {13},{BC} = {10},{PA} = {12}$ ,则 $P$ 到 ${BC}$ 的距离等于\blkda{${12}\sqrt{2}$}
        \begin{jieda}
            作 ${AD}\bot {BC}$ 于 $D$ ,连接 ${PD}$ . 不难求得 ${AD} = {12}$ . 又 ${PD}$ 在平面 ${ABC}$ 上的射影是 ${AD}$ ,根据三垂线定理, ${BC} \bot  {AD} \Longrightarrow  {BC} \bot  {PD}$ ,故 $P$ 到 ${BC}$ 的距离即为 ${PD} = \sqrt{P{A}^{2} + D{A}^{2}} = {12}\sqrt{2}$ .
        \end{jieda}
        \item 已知 ${PA}$ 垂直于正六边形 ${ABCDEF}$ 所在平面,且 ${PA} = 1$ ,边长 ${AB} = 2$ ,则点 $P$ 到边 ${BC}$ 的距离等于\blkda{2}
        \begin{jieda}
            作 ${AG} \bot  {CB}$ 的延长线于 $G$ ,连接 ${PG}$ . 不难求得 ${AG} = \sqrt{3}$ . 又 ${PG}$ 在正六边形 ${ABCDEF}$ 所在平面上的射影是 ${AG}$ ,根据三垂线定理, ${CB} \bot  {AG} \Longrightarrow  {CB} \bot  {PG}$ ,故 $P$ 到 ${BC}$ 的距离即为 ${PG} = \sqrt{P{A}^{2} + G{A}^{2}} = 2$ .
        \end{jieda}
        \item 过平面 $\alpha$ 外一点 $P$ 引 $\alpha$ 的垂线 ${PO}$ 与斜线 ${PA},{PB}(O,A,B$ 为垂足与斜足),若 ${PA} =$  $8{{cm}},{PB} = 5{{cm}}$ ,又 ${OA} : {OB} = 4 : \sqrt{3}$ ,则点 $P$ 到 $\alpha$ 的距离等于\blkda{4} ${{cm}}$
        \begin{jieda}
            设点 $P$ 到 $\alpha$ 的距离 ${PO} = h {{cm}}$ ,并设 ${OA} = {4k}{{cm}},{OB} = \sqrt{3}k {{cm}}\left( {h,k > 0}\right)$ . 在Rt $\bigtriangleup  {POA}$ 和 Rt $\bigtriangleup  {POB}$ 中,由勾股定理得 $\left\{  \begin{array}{l} {h}^{2} + {16}{k}^{2} = {64}, \\  {h}^{2} + 3{k}^{2} = {25}. \end{array}\right.$ 解得 $\left\{  \begin{array}{l} h = 4 \\  k = \sqrt{3} \end{array}\right.$ .
        \end{jieda}
    \end{enumerate}

    \item 已知矩形 ${ABCD}$ 的边长 ${AB} = 6{{cm}},{BC} = 4{{cm}}$ ,在 ${CD}$ 上截取 ${CE} =$  $4{{cm}}$ ,以 ${BE}$ 为棱将矩形折起,使 $\bigtriangleup B{C}^{\prime }E$ 的高 ${C}^{\prime }F \bot$ 平面 ${ABED}$ (如图), 求:
    
    \begin{center}
    \includegraphics[width=0.2\textwidth]{bodyxb/src/bo_d20aucref24c73anc48g_149_1145_141_264_187_0.jpg}
    \end{center}
    \begin{enumerate}
        \item 点 ${C}^{\prime }$ 到平面 ${ABED}$ 的距离
        \item 点 ${C}^{\prime }$ 到边 ${AB}$ 的距离
        \item 点 ${C}^{\prime }$ 到 ${AD}$ 的距离
    \end{enumerate}
    \begin{jieda}
        作 ${FH}\bot {AB}$ 于 $H$ ,作 ${FG}\bot {AD}$ 于 $G$ ,则 ${C}^{\prime }H\bot {AB}$ , ${C}^{\prime }G\bot {AD}$
        \begin{enumerate}
            \item ${C}^{\prime }$ 到面 ${ABED}$ 的距离为 ${C}^{\prime }F = 2\sqrt{2}{{cm}}$ .
            \item ${C}^{\prime }$ 到 ${AB}$ 的距离为 ${C}^{\prime }H = 2\sqrt{3}{{cm}}$ .
            \item ${C}^{\prime }$ 到 ${AD}$ 的距离为 ${C}^{\prime }G = 2\sqrt{6}{{cm}}$ .
        \end{enumerate}
    \end{jieda}
\end{enumerate}


\section*{线面角}
\subsection*{线面角}
\begin{enumerate}
    \item 如图，在正三棱柱 ${ABC} - {A}_{1}{B}_{1}{C}_{1}$ 中， ${AB} = 1,{A{A}_{1}} = \sqrt{2}$ . 则直线 $A{C}_{1}$ 与平面 ${AB}{B}_{1}{A}_{1}$ 所成角的大小为\blkda{${30}^{ \circ  }$}
    
    \begin{center}
    \includegraphics[width=0.2\textwidth]{bodyxb/src/bo_d4q3m9bef24c73bcm6mg_2_261_1302_308_348_0.jpg}
    \end{center}

    \begin{jieda}
        取 ${A}_{1}{B}_{1}$ 的中点 $D$ ,连接 ${C}_{1}D\text{ 、 }{AD}$ ,如图:
        
        \begin{center}
        \includegraphics[width=0.2\textwidth]{bodyxb/src/bo_d4q3m9bef24c73bcm6mg_3_254_210_315_356_0.jpg}
        \end{center}
        
        
        因为三棱柱 ${ABC} - {A}_{1}{B}_{1}{C}_{1}$ 是正三棱柱，所以 $\bigtriangleup  {A}_{1}{B}_{1}{C}_{1}$ 为等边三角形,所以 ${C}_{1}D\bot {A}_{1}{B}_{1}$ , 又 $A{A}_{1} \bot$ 平面 ${A}_{1}{B}_{1}{C}_{1},{C}_{1}D \subset$ 平面 ${A}_{1}{B}_{1}{C}_{1}$ ,所以 ${C}_{1}D \bot  A{A}_{1}$，又 ${A}_{1}{B}_{1} \cap  A{A}_{1} = {A}_{1},{A}_{1}{B}_{1},A{A}_{1} \subset$ 平面 ${AB}{B}_{1}{A}_{1}$ ,所以 ${C}_{1}D \bot$ 平面 ${AB}{B}_{1}{A}_{1}$，所以 $\angle {C}_{1}{AD}$ 为直线 $A{C}_{1}$ 与平面 ${AB}{B}_{1}{A}_{1}$ 所成角.
        
        又 ${C}_{1}D = \sqrt{{1}^{2} - {\left( \dfrac{1}{2}\right) }^{2}} = \dfrac{\sqrt{3}}{2},A{C}_{1} = \sqrt{{1}^{2} + {\left( \sqrt{2}\right) }^{2}} = \sqrt{3}$，$\sin \angle {C}_{1}{AD} = \dfrac{{C}_{1}D}{A{C}_{1}} = \dfrac{1}{2}$ ,所以 $\angle {C}_{1}{AD} = {30}^{ \circ  }$，即直线 $A{C}_{1}$ 与平面 ${AB}{B}_{1}{A}_{1}$ 所成角的大小为 ${30}^{ \circ  }$ .
    \end{jieda}
    
    \item 正四面体 $P - {ABC}$ 中,直线 ${PC}$ 与平面 ${PAB}$ 所成角的余弦值为\blkda{$\dfrac{\sqrt{3}}{3}$}

    \begin{jieda}
        如图,取 ${AB}$ 中点 $D$ ,在线段 ${PD}$ 上取点 $E$ ,使得 ${PE} = {2ED}$ ,连接 ${DP},{DC},{EC}$ .
        
        \begin{center}
        \includegraphics[width=0.2\textwidth]{bodyxb/src/bo_d4q3m9bef24c73bcm6mg_2_243_207_484_495_0.jpg}
        \end{center}
        
        
        因为 $P - {ABC}$ 为正四面体，所以 $E$ 为正 $\bigtriangleup  {PAB}$ 的重心，且 ${CE}\bot$ 平面 ${PAB}$，所以 $\angle {CPD}$ 即为直线 ${PC}$ 与平面 ${PAB}$ 所成的角.
        
        不妨设 ${AB} = 2$ ，则在 $\bigtriangleup  {PCE}$ 中， ${PE}\bot {CE}$ ， ${PE} = \dfrac{2}{3}{PD} = \dfrac{2\sqrt{3}}{3}$ ， ${PC} = 2$ ，所以 $\cos \angle {CPD} = \dfrac{PE}{PC} = \dfrac{\dfrac{2\sqrt{3}}{3}}{2} = \dfrac{\sqrt{3}}{3}$ .
    \end{jieda}
    
    \item 正三棱锥 $P - {ABC}$ 中，$AB = 1, PA = 3$，直线 ${PC}$ 与平面 ${ABC}$ 所成角的正弦值为\blkda{$\dfrac{\sqrt{78}}{9}$}

    \item 如图，在棱长为 1 的正方体 ${ABCD} - {A}_{1}{B}_{1}{C}_{1}{D}_{1}$ 中，那么直线 ${B}_{1}{D}_{1}$ 与平面 ${AC}{D}_{1}$ 所成的角的余弦值是\blkda{$\dfrac{\sqrt{3}}{3}$}

    \begin{center}
    \includegraphics[width=0.25\textwidth]{bodyxb/src/bo_d4q3m9bef24c73bcm6mg_4_241_1380_340_320_0.jpg}
    \end{center}
    \begin{jieda}
        连接$D_1$和$AC$中点$O$，取$D_1G = \dfrac{2}{3}D_1O$，易知$G$为$\triangle ACD_1$的重心，$B_1G \bot \triangle ACD_1$，故$\angle{B_1D_1G}$即为线面角.

        $D_1G = \dfrac{\sqrt{6}}{3}$，$D_1B_1 = \sqrt{2}$，故$cos\angle{B_1D_1G} = \dfrac{\sqrt{3}}{3}$
    \end{jieda}


    \item 如图,在正方体 ${ABCD} - {A}_{1}{B}_{1}{C}_{1}{D}_{1}$ 中,点 $O$ 为线段 ${BD}$ 的中点,设点 $P$ 在线段 $C{C}_{1}$ 上,直线 ${OP}$ 与平面 ${A}_{1}{BD}$ 所成的角为 $\alpha$ ,则 $\sin \alpha$ 的最小值为\blkda{$\dfrac{\sqrt{6}}{3}$}
    
    \begin{center}
    \includegraphics[width=0.2\textwidth]{bodyxb/src/bo_d4ojskjef24c73bc0dj0_2_245_749_414_371_0.jpg}
    \end{center}

    \begin{jieda}
        易知$A_1, O, P \in $面$AA_1C_1C$，做$PH \bot A_1O$，故$H \in A_1O$，故$\angle{POA_1}$即为线面所成角或其补角.

        注意到 $A{C}_{1} \bot$ 平面 ${A}_{1}{BD}$，所以移动 $P$ 的位置，当$OP$和$AC_1$平行时取得最大值，故当 $P$ 在 $C$ 或 ${C}_{1}$ 处时,线面角 $\alpha$ 才可能最小.
    
        在平面$AA_1C_1C$中，计算可得$\sin \alpha$ 的最小值为 $\dfrac{\sqrt{6}}{3}$ .
    \end{jieda}
    
\end{enumerate}

\subsection*{三余弦定理与最小角定理}
\begin{enumerate}
    \item ${PA},{PB},{PC}$ 是从点 $P$ 出发的不共面的三条射线,每两条射线的夹角均为 ${60}^{ \circ  }$ ,那么直线 ${PC}$ 与平面 ${PAB}$ 所成角的余弦值是\blkda{$\dfrac{\sqrt{3}}{3}$}
    \begin{jieda}
        如图,设直线 ${PC}$ 在平面 ${PAB}$ 的射影为 ${PD}$ ,
        
        \begin{center}
        \includegraphics[width=0.2\textwidth]{bodyxb/src/bo_d4ojskjef24c73bc0dj0_1_243_274_288_325_0.jpg}
        \end{center}

        根据三余弦定理可知$cos\angle{APC} = cos\angle{APD} \cdot cos\angle{CPD}$

        易知$\angle{APC} = \ang{60}$，$\angle{APD} = \ang{30}$，故$cos\angle{CPD} = \dfrac{\sqrt{3}}{3}$.
    \end{jieda}

    \item 已知三棱柱 ${ABC} - {A}_{1}{B}_{1}{C}_{1}$ 的侧棱与底面边长都相等， ${A}_{1}$ 在底面 ${ABC}$ 上的射影为 ${BC}$ 的中点,则异面直线 ${AB}$ 与 $C{C}_{1}$ 所成的角的余弦值为\blkda{$\dfrac{3}{4}$}
    \begin{jieda}
        设 ${BC}$ 的中点为 $D$，由题意可知 ${A}_{1}D \bot$ 平面 ${ABC}$，连接 ${A}_{1}D\text{ 、 }{AD}\text{ 、 }{A}_{1}B$ ,易知 $\theta  = \angle {A}_{1}{AB}$ 即为异面直线 ${AB}$ 与 $C{C}_{1}$ 所成的角.

        \begin{center}
        \includegraphics[width=0.2\textwidth]{bodyxb/src//bo_d4op0rv7aajc73fs5h4g_0_265_1176_414_277_0.jpg}
        \end{center}

        设三棱柱 ${ABC} - {A}_{1}{B}_{1}{C}_{1}$ 的侧棱与底面边长为 2，根据几何关系可知$A_1D = 1$，故$\angle{A_1AD} = \angle{BAD} = \ang{30}$，故根据三余弦定理可知$cos\angle{A_1AB} = \dfrac{3}{4}$.
    \end{jieda}
    \item 已知三棱柱 ${ABC} - {A}_{1}{B}_{1}{C}_{1}$ 的侧棱与底面边长都相等， ${A}_{1}$ 在底面 ${ABC}$ 上的射影为底面的中心，则 ${AB_1}$ 与 底面$ABC$ 所成的角的正弦值为\blkda{$\dfrac{\sqrt{2}}{3}$}
    \begin{jieda}
        设底面的中心为 $O$，由题意可知 ${A}_{1}O \bot$ 平面 ${ABC}$
        
        设三棱柱 ${ABC} - {A}_{1}{B}_{1}{C}_{1}$ 的侧棱与底面边长为 1，根据几何关系可知$AO = \dfrac{\sqrt{3}}{3}$，故$A_1O = \dfrac{\sqrt{6}}{3}$，故$cos\angle{A_1AO} = \dfrac{\sqrt{3}}{3}$，又因为$\angle{BAO} = \dfrac{\pi}{6}$，故根据三余弦定理可知$cos\angle{A_1AB} = \dfrac{1}{2}$.

        因此$\angle{A_1AB} = \ang{60}$，故$AB_1 = 2\sqrt{3}$，所以所求角的正弦值为$\dfrac{B_1O_1}{AB_1} = \dfrac{A_1O_1}{AB_1} = \dfrac{\sqrt{2}}{3}$
    \end{jieda}
    \item 已知 $\angle {ACB} = {90}^{ \circ  }$ ， $P$ 为平面 ${ABC}$ 外一点， ${PC} = 2$ ，点 $P$ 到 $\angle {ACB}$ 两边 ${AC}$ ， ${BC}$ 的距离均为 $\sqrt{3}$ ，那么 $P$ 到平面 ${ABC}$ 的距离为\blkda{$\sqrt{2}$}.
    \begin{jieda}

        做$PO \bot $面$ACB$，根据对称性可知$\angle{ACO} = \angle{BCO} = \ang{45}$，根据几何关系可知$\angle{PCB} = \angle{PCA} = \ang{60}$，故根据三余弦定理可知$\angle{PCO} = \ang{45}$，所以$PO = CO = \sqrt{2}$
        
        \begin{center}
        \includegraphics[width=0.2\textwidth]{bodyxb/src//bo_d4op0rv7aajc73fs5h4g_1_245_1319_416_484_0.jpg}
        \end{center}
    \end{jieda}


    \item 如图,已知正方体 ${ABCD} - {A}_{1}{B}_{1}{C}_{1}{D}_{1}$ 的棱长为 $2,M,N$ 分别为棱 ${AB},{BC}$ 的中点,若点 $P$ 为线段 ${B}_{1}N$ 上的动点,不包括端点),设异面直线 ${C}_{1}P$ 与 ${MN}$ 所成的角为 $\theta$ ,则 $\cos \theta$ 的取值范围是\blkda[2]{$\left( {\dfrac{\sqrt{10}}{10},\dfrac{\sqrt{2}}{2}}\right)$}
    
    \begin{center}
    \includegraphics[width=0.2\textwidth]{bodyxb/src/bo_d4ojskjef24c73bc0dj0_5_239_263_376_348_0.jpg}
    \end{center}
    \begin{jieda}
        如图所示,连接 ${AC},{A}_{1}{C}_{1}$ ,因为 $M,N$ 分别为棱 ${AB},{BC}$ 的中点,可得 ${MN}//{AC}$ 在正方体 ${ABCD} - {A}_{1}{B}_{1}{C}_{1}{D}_{1}$ 中,可得 ${A}_{1}{C}_{1}//{AC}$ ,所以 ${MN}//{A}_{1}{C}_{1}$，所以$\angle {A}_{1}{C}_{1}P$为异面直线 ${C}_{1}P$ 与 ${MN}$ 所成的角或其补角.

        由三余弦定理可知$cos\angle{A_1C_1P} = cos\angle{A_1C_1B_1} \cdot cos\angle{PC_1B_1} = \dfrac{\sqrt{2}}{2} \cdot cos\angle{PC_1B_1}$

        因为$cos\angle{PC_1B_1} \in \left(\dfrac{\sqrt{5}}{5}, 1\right)$，所以$cos\angle{A_1C_1P} \in \left( {\dfrac{\sqrt{10}}{10},\dfrac{\sqrt{2}}{2}}\right)$
    \end{jieda}
    
    \item 如图所示，在侧棱垂直于底面的三棱柱 ${ABC} - {A}_{1}{B}_{1}{C}_{1}$ 中， $P$ 是棱 ${BC}$ 上的动点，记直线 ${A}_{1}P$ 与平面 ${ABC}$ 所成的角为 ${\theta }_{1}$ ,与直线 ${BC}$ 所成的角为 ${\theta }_{2}$ ,则 ${\theta }_{1},{\theta }_{2}$ 的大小关系是
    
    \begin{center}
    \includegraphics[width=0.2\textwidth]{bodyxb/src/bo_d4ojskjef24c73bc0dj0_0_242_626_445_446_0.jpg}
    \end{center}

    \xx{${\theta }_{1} = {\theta }_{2}$}{${\theta }_{1} > {\theta }_{2}$}{${\theta }_{1} < {\theta }_{2}$}{不能确定}

    \begin{jieda}
        根据线面角是该直线与对应平面内的任意直线所成角中最小的角，所以有 ${\theta }_{1} < {\theta }_{2}$ ,故选 C.
    \end{jieda}

    \item 在侧棱垂直于底面的三棱柱 ${ABC} - {A}_{1}{B}_{1}{C}_{1}$ 中，$\angle{ABC} = \ang{90}$，直线$l$和面$ABB_1A_1$和面$ACC_1A_1$的夹角为$\theta_1, \theta_2$,则 ${\theta }_{1},{\theta }_{2}$ 满足\dotfill{(\kh{C})}
    \xx{${\theta }_{1}+{\theta }_{2} = \dfrac{\pi}{2}$}{${\theta }_{1}+{\theta }_{2} \geq \dfrac{\pi}{2}$}{${\theta }_{1}+{\theta }_{2} \leq \dfrac{\pi}{2}$}{${\theta }_{1}+{\theta }_{2} < \dfrac{\pi}{2}$}

    \begin{jieda}
        设直线$l$和面$ABB_1A_1$和面$ACC_1A_1$交于$P, Q$，过$P, Q$分别$AA_1$的垂线，设垂足为$H_1, H_2$，则$\theta_1 = \angle{QPH_2}, \theta_2 = \angle{PQH_1}$，故$\theta_1 \leq \angle{PQH_2}$，$\theta_2 + \angle{PQH_2} = \ang{90}$
    \end{jieda}

    \item 已知等腰$Rt \bigtriangleup {ABC}$ 内接于圆 $O$，连接$OC$，交点圆 $O$于$C'$， $M$ 是下半圆弧上的动点 (如图所示),现将上半圆面沿 ${AB}$ 折起，使得$\angle{COC'} = \dfrac{\pi }{4}$ . 则直线 ${AC}$ 与直线 ${OM}$ 所成角的最小值是\blkda{$\dfrac{\pi }{6}$}
    
    \begin{center}
    \includegraphics[width=0.2\textwidth]{bodyxb/src/bo_d4onr53ef24c73bc3ut0_0_246_691_285_267_0.jpg}
    \end{center}
    \begin{jieda}
        ${AC}$ 与直线 ${OM}$ 所成角的最小值是$AC$与平面$ABM$的线面所成角

        做$CH \bot $面$AOM$，则$CH = \dfrac{r}{\sqrt{2}}$，$AC = \sqrt{2}r$，故线面所成角为$\dfrac{\pi}{6}$.
    \end{jieda}
    
    \item 如图,已知正三棱锥 ${A-BCD},{BC} = {CD} = {BD} = \sqrt{3},{AB} = {AC} = {AD} = 2$ ,点 $P,Q$ 分别在棱 ${BC},{CD}$ 上(不包含端点)，则直线 ${AP},{BQ}$ 所成的角的取值范围是\blkda{$\left( {\dfrac{\pi }{3},\dfrac{\pi }{2}}\right\rbrack$}
    
    \begin{center}
    \includegraphics[width=0.2\textwidth]{bodyxb/src/bo_d4ojskjef24c73bc0dj0_6_241_910_382_394_0.jpg}
    \end{center}

    \begin{jieda}
        如下图所示:
        
        \begin{center}
        \includegraphics[width=0.2\textwidth]{bodyxb/src/bo_d4ojskjef24c73bc0dj0_6_242_1650_380_393_0.jpg}
        \end{center}

        当$P, Q$分别为中点时$AP \bot BQ$

        根据最小角定理可知$AP$和$BQ$所成角大于等于$AP$和面$BCD$的夹角，$\angle{AOP}$即为线面角，而$P$位于$BC$端点处会有$\angle{APO}$的下限$\dfrac{\pi}{3}$，故${AP},{BQ}$ 所成的角的取值范围是$\left( {\dfrac{\pi }{3},\dfrac{\pi }{2}}\right\rbrack$.
    \end{jieda}
\end{enumerate}

\section*{面面平行}
\section*{面面垂直}
\section*{二面角}
\subsection*{三正弦定理与最大角定理}
\section*{证明专题}
\subsection*{证明线线平行}
证明线线平行涉及到的方法如下：
\begin{dingautolist}{172}
    \item 平面几何：中位线、平行四边形
    \item 平行公理（公理4）：平行的传递性
    \item 线面平行的性质定理：线面平行中蕴含着线线平行
    \item 线面垂直的性质定理：垂直和平行可以相互转化
    \item 面面平行的性质定理：面面平行中蕴含着线线平行
\end{dingautolist}
\subsubsection*{平行公理}
\begin{enumerate}
    \item 如图,正方体的棱长是 $a$ ,点 $E\text{ 、 }F$ 分别是两条棱的中点. 求证:四边形 ${BDEF}$ (图中阴影部分)是一个梯形
    
    \begin{center}
    \includegraphics[width=0.2\textwidth]{bodyxb/src/bo_d4lnl5ref24c739do320_11_241_208_336_326_0.jpg}
    \end{center}
    \begin{jieda}
        如图所示，连接 ${B}_{1}{D}_{1}$
        
        \begin{center}
        \includegraphics[width=0.2\textwidth]{bodyxb/src/bo_d4lnl5ref24c739do320_11_243_1097_335_324_0.jpg}
        \end{center}
        
        
        因为点 $E$ 、 $F$ 分别是 ${A}_{1}{D}_{1}$ 、 ${A}_{1}{B}_{1}$ 的中点，所以 ${EF}//{B}_{1}{D}_{1}$ ， ${EF} = \dfrac{1}{2}{B}_{1}{D}_{1}$ ，
        
        又因为 ${BD}//{B}_{1}{D}_{1},{BD} = {B}_{1}{D}_{1}$ ,所以 ${EF}//{BD},{EF} = \dfrac{1}{2}{BD}$ ,
        
        所以四边形 ${BDEF}$ 是一个梯形.
    \end{jieda}
    \item 如图所示，在空间四边形 ${ABCD}$ 中，点 $E$ ， $H$ 分别是边 ${AB}$ ， ${AD}$ 的中点，点 $F$ ， $G$ 分别是边 ${BC},{CD}$ 上的点，且 $\dfrac{CF}{CB} = \dfrac{CG}{CD} = \dfrac{2}{3}$ ，有以下结论正确的是\blkda{\ding{173}\ding{175}}
    
    \begin{center}
    \includegraphics[width=0.2\textwidth]{bodyxb/src/bo_d4lnl5ref24c739do320_2_243_995_278_307_0.jpg}
    \end{center}

    \begin{dingautolist}{172}
        \item $EF$与$GH$平行
        \item $EF$与$GH$共面
        \item $EF$与$GH$的交点 $M$ 可能在直线 ${AC}$ 上，也可能不在直线 ${AC}$ 上
        \item ${EF}$ 与 ${GH}$ 的交点 $M$ 一定在直线 ${AC}$ 上.
    \end{dingautolist}

    \begin{jieda}
        如图所示. 连接 ${EH},{FG}$ ,
        
        \begin{center}
        \includegraphics[width=0.2\textwidth]{bodyxb/src/bo_d4lnl5ref24c739do320_3_239_193_280_308_0.jpg}
        \end{center}
        
        
        依题意,可得 ${EH}//{BD},{FG}//{BD}$，所以 ${EH}//{FG}$，所以 $E,F,G,H$ 共面
        
        因为 ${HE} = \dfrac{1}{2}{BD},{FG} = \dfrac{2}{3}{BD}$，所以四边形 ${EFGH}$ 是梯形, ${EF}$ 与 ${GH}$ 必相交，设EF与 ${GH}$ 的交点为 $M$
        
        因为点 $M$ 在 ${EF}$ 上，故点 $M$ 在平面 ${ACB}$ 上；同理，点 $M$ 在平面 ${ACD}$ 上,
        
        所以点 $M$ 在平面 ${ACB}$ 与平面 ${ACD}$ 的交线上， 又 ${AC}$ 是这两个平面的交线，所以点 $M$ 一定在直线 ${AC}$ 上
    \end{jieda}
    \item 如图,在三棱柱 ${ABC} - {{A}_{1}{B}_{1}{C}_{1}}$ 中， $E,F,G,H$ 分别为 $B{B}_{1},C{C}_{1},{A}_{1}{B}_{1},{A}_{1}{C}_{1}$ 的中点,则下列说法错误的是\dotfill{(\kh{D})}
    
    \begin{center}
    \includegraphics[width=0.2\textwidth]{bodyxb//src/bo_d4lsgh3ef24c739dpshg_0_240_621_292_321_0.jpg}
    \end{center}

    \xx{$E,F,G,H$ 四点共面}{${EF}//{GH}$}{${EG},{FH},A{A}_{1}$ 三线共点}{$\angle {EG}{B}_{1} = \angle {FH}{C}_{1}$}
    \begin{jieda}
        对于 ${{AB}}$ ,如图,连接 ${EF},{GH}$，因为 ${GH}$ 是 $\bigtriangleup {A}_{1}{B}_{1}{C}_{1}$ 的中位线，所以 ${GH}//{B}_{1}{C}_{1}$
        
        因为 ${B}_{1}E//{C}_{1}F$ ,且 ${B}_{1}E = {C}_{1}F$ ,所以四边形 ${B}_{1}{EF}{C}_{1}$ 是平行四边形，所以 ${EF}//{{B}_{1}{C}_{1}}$ ，所以 ${EF}//{GH}$ ，所以 $E,F,G,{HH}$ 四点共面，故 ${{AB}}$ 正确
        
        \begin{center}
        \includegraphics[width=0.2\textwidth]{bodyxb//src/bo_d4lsgh3ef24c739dpshg_0_241_1397_296_419_0.jpg}
        \end{center}
        
        对于 ${C}$ ,如图,延长 ${EG},{FH}$ 相交于点 $P$ ,
        
        因为 $P \in  {EG}$ , ${EG} \subset$ 平面 ${AB}{B}_{1}{A}_{1}$ ,所以 $P \in$ 平面 ${AB}{B}_{1}{A}_{1}$
        
        因为 $P \in  {FH},{FH} \subset$ 平面 ${AC}{C}_{1}{A}_{1}$ ,所以 $P \in$ 平面 ${AC}{C}_{1}{A}_{1}$
        
        因为平面 ${AB}{B}_{1}{A}_{1} \cap$ 平面 ${AC}{C}_{1}{A}_{1} = A{A}_{1}$，所以 $P \in  A{A}_{1}$ ,所以 ${EG},{FH},A{A}_{1}$ 三线共点,故 ${C}$ 正确;
        
        对于 ${D}$ ,因为 $E{B}_{1} = F{C}_{1}$ ,当 $G{B}_{1} \neq  H{C}_{1}$ 时, $\tan \angle {EG}{B}_{1} \neq  \tan \angle {FH}{C}_{1}$ ,
        
        又 $0 < \angle {EG}{B}_{1},\angle {FH}{C}_{1} < \dfrac{\pi }{2}$ ,则 $\angle {EG}{B}_{1} \neq  \angle {FH}{C}_{1}$ ,故 ${D}$ 错误.
    \end{jieda}
\end{enumerate}
\subsubsection*{线面平行性质定理}
\begin{enumerate}
    \item 如图所示，在四棱锥 $P - {ABCD}$ 中， $M$ ， $N$ 分别为 ${PC}$ ， ${AC}$ 上的点，且 ${MN}//$ 平面 ${PAD}$ ，则下列说法正确的是\dotfill{(\kh{B})}

    \begin{center}
    \includegraphics[width=0.2\textwidth]{bodyxb/src/bo_d4lnl5ref24c739do320_0_242_545_335_301_0.jpg}
    \end{center}

    \xx{$MN//PD$}{$MN//PA$}{${MN}//{AD}$}{以上均有可能}
    \begin{jieda}
        直线 ${MN} \subset$ 平面 ${PAC},{MN}//$ 平面 ${PAD}$ ,平面 ${PAC} \cap$ 平面 ${PAD} = {PA}$，所以$MN//PA$.
    \end{jieda}

    \item 求证：已知三个平面两两相交，且有三条交线. 如果这三条交线中有两条交线平行，那么第三条交线也与另两条直线平行.
    \begin{jieda}
        设$\alpha \cap \beta = a, \beta \cap \gamma = b, \alpha \cap \gamma = c$

        $a // b$则根据线面平行的判定定理可知$a // \beta$，故根据线面平行性质定理可知$a // c$，根据公理4可知$b //c$.
    \end{jieda}
    
    \item 求证: 如果一条直线与两个相交平面都平行, 那么这条直线与这两个平面的交线平行.

    即：已知直线 $l//$ 平面 $\alpha ,l//$ 平面 $\beta$ ,且 $\alpha  \cap  \beta  = m$ (如图),求证: $l//m$ .
    \begin{center}
    \includegraphics[width=0.2\textwidth]{bodyxb/src/bo_d20aucref24c73anc48g_141_342_212_217_226_0.jpg}
    \end{center}
    \begin{jieda}
        \begin{enumerate}
            \item 方法1
            
            过 $l$ 作平面 $\gamma$ ,使 $\gamma  \cap  \alpha  = a$ ,过 $l$ 作平面 $\delta$ ,使 $\delta  \cap  \beta  = b$ ,则 $l//a,l//b$ ,于是 $a//b$ ,故 $a//\beta$ ,由此得 $a//m,\therefore l//m$

            \item 方法2
            
            在$m$上取点$A$，则根据公理2的推论可知点$A$和$l$确定一个平面，记为$\gamma$，设$\gamma \cap \alpha = l_1, \gamma \cap \beta = l_2$，则根据线面平行性质定理可知$l // l_1, l // l_2$

            因为在$\gamma$内过一点$A$有且仅有一条直线与$l$平行，所以$l_1 = l_2$
            
            因此$\gamma \cap \beta = l_1$，所以$l_1 \subset \alpha, l_1 \subset \beta$，故根据公理3可知$l_1 = m$
    
            故$m // l$
        \end{enumerate}
    \end{jieda}

    \item 如图所示，在多面体 ${A}_{1}{B}_{1}{D}_{1}-{DCBA}$ ，四边形 $A{A}_{1}{B}_{1}B$ ， ${AD}{D}_{1}{A}_{1}$ ， ${ABCD}$ 均为正方形， $E$ 为 ${B}_{1}{D}_{1}$ 的中点,过 ${A}_{1},D,E$ 的平面交 $C{D}_{1}$ 于 ${F}$，证明: ${EF}//{B}_{1}C$
    
    \begin{center}
    \includegraphics[width=0.2\textwidth]{bodyxb/src/bo_d4lnl5ref24c739do320_5_243_863_315_272_0.jpg}
    \end{center}

    \begin{jieda}
        依据正方形的性质可知 ${A}_{1}{B}_{1}//{AB}//{DC}$ ,且
        
        ${A}_{1}{B}_{1} = {AB} = {DC}$ ,从而 ${A}_{1}{B}_{1}{CD}$ 为平行四边形,则 ${B}_{1}C//{A}_{1}D$ ,根据线面平行的判定定理知 ${B}_{1}C//$ 面 ${A}_{1}{DE}$ ，再由线面平行的性质定理知 ${EF}//{B}_{1}C$ .
    \end{jieda}

    \item 如图, 四棱锥P-ABCD中, ABCD是平行四边形, $M$是$PC$的中点.在 ${DM}$ 上取一点 $G$ ，过 $G$ 和 ${AP}$ 作平面交平面 ${BDM}$ 于 ${GH}$ ， $H$ 在 ${BD}$ 上，证明: ${AP}//{GH}$ .
    
    \begin{center}
    \includegraphics[width=0.2\textwidth]{bodyxb/src/bo_d4lnl5ref24c739do320_6_242_1523_364_338_0.jpg}
    \end{center}

    \begin{jieda}
        连接 ${AC}$ 与 ${BD}$ 交于点 $O$ ，且 $O$ 为 ${AC}$ 的中点，
        
        \begin{center}
        \includegraphics[width=0.2\textwidth]{bodyxb/src/bo_d4lnl5ref24c739do320_7_239_1373_368_337_0.jpg}
        \end{center}
        
        由点 $M$ 为 ${PC}$ 的中点,所以 ${AP}//{MO}$ ,
        
        因为 ${AP} \not \subset$ 平面 ${BDM}$ ,且 ${MO} \subset$ 平面 ${BDM}$ ,所以 ${AP}//$ 平面 ${BDM}$ ,
        
        又因为 ${AP} \subset$ 平面 ${AHGP}$ ,且平面 ${AHGP} \cap$ 平面 ${BDM} = {GH}$ ,所以 ${AP}//{GH}$ .
        
    \end{jieda}
    
     \item 如图，在四棱锥 $P - {ABCD}$ 中，底面 ${ABCD}$ 是边长为 $a$ 的正方形， $O$ 为正方形 ${ABCD}$ 的中心， ${PO} \bot$ 平面 ${ABCD}$ . 若点 $M$ 在棱 ${PC}$ 上且不与 $P$ 、 $C$ 重合，平面 ${ADM}$ 交棱 ${PB}$ 于点 $N$ ，求证: ${MN}//{AD}$ .
    
    \begin{center}
    \includegraphics[width=0.2\textwidth]{bodyxb/src/bo_d4lnl5ref24c739do320_12_241_500_373_340_0.jpg}
    \end{center}
    \begin{jieda}
        因为${AD}//{BC}$ 且 ${AD} \subset$ 平面 ${ADM}$ ， ${BC}$ 不在平面 ${ADM}$ 上， 所以${BC}//$ 平面 ${ADM}$ ，

        又 因为${BC} \subset$ 平面 ${PBC}$ ,平面 ${ADM} \cap$ 平面 ${PBC} = {MN}$，因此$ {BC}//{MN}$ ,且 ${BC}//{AD}$，所以${MN}//{AD}$ .
    \end{jieda}
    
    \item 如图所示，已知三棱柱 ${ABC} - {A}^{\prime }{B}^{\prime }{C}^{\prime }$ 中，$D$ 是 $BC$ 的中点，$D'$是 $B^{\prime }{C}^{\prime }$ 的中点，设平面 ${A}^{\prime }{D}^{\prime }B \cap$ 平面 ${{ABC}} = {a}$ ,平面 ${{AD}}{{C}}^{\prime } \cap$ 平面 ${{A}}^{\prime }{{B}}^{\prime }{{C}}^{\prime } = {b}$ ,判断直线 ${a},{b}$ 的位置关系,并证明.
    
    \begin{center}
    \includegraphics[width=0.2\textwidth]{bodyxb/src/bo_d4lnl5ref24c739do320_4_265_1082_306_354_0.jpg}
    \end{center}
    \begin{jieda}
        连接$DD'$，因为$D$ 是 $BC$ 的中点，$D'$是 $B{C}^{\prime }$ 的中点，所以$DD' // BB'$，且$DD' = BB'$，故$DD' // AA'$，且$DD' = AA'$，故$A'D' // AD$，且$A'D' = AD$，故$AD // $平面$A'B'C'$，$A'D' // $平面$ABC$
    
        因为平面 ${{A}}^{\prime }{{D}}^{\prime }{B} \cap$ 平面 ${{ABC}} = {a}$，$A'D' // $平面$ABC$，所以${{A}}^{\prime }{{D}}^{\prime }//{a}$，同理可得 ${{AD}}//{b}$，根据平行的传递性可知 ${a}//{b}$ .
    \end{jieda}
\end{enumerate}
\subsubsection*{线面垂直性质定理}
\begin{enumerate}
    \item 若 $a,b$ 表示两条直线， $\alpha$ 表示一个平面，则下列命题为真的是\blkda{\ding{173}\ding{174}}
    \begin{dingautolist}{172}
        \item 若 $a//b,b \subset  \alpha$ ,则 $a//\alpha$
        \item 若 $a \bot  \alpha ,a//b$ ,则 $b \bot  \alpha$
        \item 若 $a \bot  \alpha ,b \bot  \alpha$ ,则 $a//b$
        \item 若 $a//\alpha ,b \subset  \alpha$ ,则 $a//b$
    \end{dingautolist}
    \begin{jieda}
        \begin{dingautolist}{172}
            \item 若 $a//b,b \subset  \alpha$ ,则 $a//\alpha$ 或 $a \subset  \alpha$ ,故错误;
            \item 若 $a \bot  \alpha ,a//b$ ,则 $b \bot  \alpha$ ,故正确;
            \item 若 $a \bot  \alpha ,b \bot  \alpha$ ,则 $a//b$ ,故正确;
            \item 若 $a//\alpha ,b \subset  \alpha$ ,则 $a$ 与 $b$ 可能平行或异面，故错误
        \end{dingautolist}
    \end{jieda}
    \item 在五面体 ${ABCDEF}$ 中, ${CD} \bot$ 平面 ${ADE},{EF} \bot$ 平面 ${ADE}$，求证: ${EF}//{AB}$
    
    \begin{center}
    \includegraphics[width=0.2\textwidth]{bodyxb/src/bo_d4lnl5ref24c739do320_9_243_1790_361_191_0.jpg}
    \end{center}
    \begin{jieda}
        因为 ${CD} \bot$ 平面 ${ADE}$ ， ${EF} \bot$ 平面 ${ADE}$ ，则 ${CD}//{EF}$ ，
        
        因为 ${EF} \not \subset$ 平面 ${ABCD},{CD} \subset$ 平面 ${ABCD}$ ,所以, ${EF}//$ 平面 ${ABCD}$ ,
        
        因为 ${EF} \subset$ 平面 ${ABFE}$ ,平面 ${ABFE} \cap$ 平面 ${ABCD} = {AB}$ ,所以, ${EF}//{AB}$ .
    \end{jieda}
    \item 如图所示,在正方体 ${ABCD} - {A}_{1}{B}_{1}{C}_{1}{D}_{1}$ 中, ${EF}$ 与 ${AC},{A}_{1}D$ 都垂直相交,垂足分别是点 $F$ 、点 $E$ .

    \begin{center}
    \includegraphics[width=0.2\textwidth]{bodyxb//src/bo_d4lsgh3ef24c739dpshg_4_242_411_406_355_0.jpg}
    \end{center}
    \begin{enumerate}
        \item 求证: ${EF}\bot$ 平面 $A{B}_{1}C$
        \item 求证: ${EF}//{B{D}_{1}}$
    \end{enumerate}
    \begin{jieda}
        \begin{enumerate}
            \item 在正方体 ${ABCD} - {A}_{1}{B}_{1}{C}_{1}{D}_{1}$ 中， ${CD}//{A}_{1}{B}_{1}$ 且 ${CD} = {A}_{1}{B}_{1}$ ，
            
            所以四边形 ${A}_{1}{B}_{1}{CD}$ 是平行四边形,所以 ${A}_{1}D//{B}_{1}C$ ,
            
            因为 ${EF} \bot  {A}_{1}D$ ,所以 ${EF} \bot  {B}_{1}C$ ,
            
            因为 ${EF} \bot  {AC},{AC} \cap  {B}_{1}C = C,{B}_{1}C \subset$ 平面 $A{B}_{1}C,{AC} \subset$ 平面 $A{B}_{1}C$ .
            
            所以 ${EF} \bot$ 平面 $A{B}_{1}C$ .
            
            \begin{center}
            \includegraphics[width=0.2\textwidth]{bodyxb//src/bo_d4lsgh3ef24c739dpshg_4_244_1688_403_357_0.jpg}
            \end{center}
            \item 因为 $D{D}_{1} \bot$ 平面 ${ABCD}$ ， ${AC} \subset$ 平面 ${ABCD}$，所以 $D{D}_{1} \bot  {AC}$ ,
            
            又因为 ${BD} \bot  {AC},D{D}_{1} \cap  {BD} = D,{BD} \subset$ 平面 ${BD}{D}_{1}{B}_{1},D{D}_{1} \subset$ 平面 ${BD}{D}_{1}{B}_{1}$ , 所以 ${AC} \bot$ 平面 ${BD}{D}_{1}{B}_{1}$ ,
            
            因为 $B{D}_{1} \subset$ 平面 ${BD}{D}_{1}{B}_{1}$ ,所以 ${AC} \bot  B{D}_{1}$ ,
            
            同理可证 $B{D}_{1} \bot  {B}_{1}C$ ,
            
            又 ${AC} \cap  {B}_{1}C = C,{B}_{1}C \subset$ 平面 $A{B}_{1}C,{AC} \subset$ 平面 $A{B}_{1}C$ .
            
            所以 $B{D}_{1} \bot$ 平面 $A{B}_{1}C$ .
            
            又 ${EF} \bot$ 平面 $A{B}_{1}C$ ,所以 ${EF}//B{D}_{1}$ .
        \end{enumerate}
    \end{jieda}

\end{enumerate}
\subsection*{证明线面平行}
\begin{enumerate}
    \item 如图，EA和DC都垂直于平面 ${ABC},{\text{ 且 }{EA} = {2DC}},{CA} = {CB}$ ，点 $F$ 为 ${BE}$ 中点，求证：$DF // $平面 ${ABC}$

    \begin{center}
    \includegraphics[width=0.2\textwidth]{bodyxb/src/bo_d4lnl5ref24c739do320_0_241_1495_252_349_0.jpg}
    \end{center}
    \begin{jieda}
        因 ${EA}$ 和 ${DC}$ 都垂直于平面 ${ABC}$ ,则 ${EA}//{DC}$
        
        \begin{center}
        \includegraphics[width=0.2\textwidth]{bodyxb/src/bo_d4lnl5ref24c739do320_1_242_1608_249_343_0.jpg}
        \end{center}
        
        取线段 ${AB}$ 的中点 $G$，连接 ${CG},{FG}$，因 $F$ 为 ${BE}$ 中点,则 ${FG}//{AE},{FG} = \dfrac{1}{2}{AE}$
        
        又 ${EA} = {2DC}$ ,则 ${CD}//{AE},{CD} = \dfrac{1}{2}{AE}$，根据平行的传递性可知${FG}//{CD}$，且${FG} = {CD}$，则四边形 ${FGCD}$ 为平行四边形,则 ${DF}//{CG}$
        
        又 ${DF} \not \subset$ 平面 ${ABC},{CG} \subset$ 平面 ${ABC}$ ,则 ${DF}//$ 平面 ${ABC}$
    \end{jieda}
    \item 如图所示,在四棱锥 $P - {ABCD}$ 中, ${BC}//$ 平面 ${PAD},{BC} = \dfrac{1}{2}{AD},E$ 是 ${PD}$ 的中点.
    
    \begin{center}
    \includegraphics[width=0.2\textwidth]{bodyxb//src/bo_d4lsgh3ef24c739dpshg_2_240_1723_425_416_0.jpg}
    \end{center}
    \begin{enumerate}
        \item 求证: ${BC}//{AD}$
        \item 求证: ${CE}//$ 平面 ${PAB}$
        \item *若 $M$ 是线段 ${CE}$ 上一动点，则线段 ${AD}$ 上是否存在点 $N$ ，使 ${MN}//$ 平面 ${PAB}$ ？说明理由
    \end{enumerate}
    \begin{jieda}
        \begin{enumerate}
            \item 在四棱锥 $P - {ABCD}$ 中， ${BC}//$ 平面 ${PAD}$ ， ${BC} \subset$ 平面 ${ABCD}$ ， ${AD} \subset$ 平面 ${PAD}$，平面 ${ABCD} \cap$ 平面 ${PAD} = {AD}$ ,所以 ${BC}//{AD}$
            
            \item 如下图，取 $F$ 为 ${AP}$ 中点，连接 ${EF},{BF}$ ，由 $E$ 是 ${PD}$ 的中点，

            所以 ${EF}//{AD}$ 且 ${EF} = \dfrac{1}{2}{AD}$ ,由上一问知 ${BC}//{AD}$ ,又 ${BC} = \dfrac{1}{2}{AD}$ ,
            
            所以 ${EF}//{BC}$ 且 ${EF} = {BC}$ ,所以四边形 ${BCEF}$ 为平行四边形,故 ${CE}//{BF}$ ,
            
            而 ${CE} \subset$ 平面 ${PAB},{BF} \not \subset $ 平面 ${PAB}$ ,则 ${CE}//$ 平面 ${PAB}$ .
            
            \begin{center}
            \includegraphics[width=0.2\textwidth]{bodyxb//src/bo_d4lsgh3ef24c739dpshg_3_243_1327_434_427_0.jpg}
            \end{center}

            \item 取 ${AD}$ 中点 $N$ ，连接 ${CN}$ ， ${EN}$ ，
            
            因为 $E,N$ 分别为 ${PD},{AD}$ 的中点,所以 ${EN}//{PA}$ ,
            
            因为 ${EN} \not \subset$ 平面 ${PAB},{PA} \subset$ 平面 ${PAB}$ ,所以 ${EN}//$ 平面 ${PAB}$ ,
            
            线段 ${AD}$ 存在点 $N$ ，使得 ${MN}//$ 平面 ${PAB}$ ，理由如下:
            
            由上一问知: ${CE}//$ 平面 ${PAB}$ ，又 ${CE} \cap  {EN} = E$ ， ${CE} \subset$ 平面 ${CEN}$ ， ${EN} \subset$ 平面 ${CEN}$ ， 所以平面 ${CEN}//$ 平面 ${PAB}$ ,又 $M$ 是 ${CE}$ 上的动点, ${MN} \subset$ 平面 ${CEN}$ ,
            
            所以 ${MN}//$ 平面 ${PAB}$ ,所以线段 ${AD}$ 存在点 $N$ ,使得 ${MN}//$ 平面 ${PAB}$ .
        \end{enumerate}
    \end{jieda}

    \item 如图,四棱锥 $P - {ABCD}$ 中, ${PA} \bot$ 底面 ${ABCD},{PA} = {AC} = 2,{BC} = 1,{AB} = \sqrt{3}$. 若 ${AD}\bot {PB}$ ，证明: ${AD}//$ 平面 ${PBC}$
    
    \begin{center}
    \includegraphics[width=0.2\textwidth]{bodyxb//src/bo_d4lsgh3ef24c739dpshg_1_242_1012_331_407_0.jpg}
    \end{center}
    \begin{jieda}
        $P$在面$ABCD$的射影是$A$，$AD \bot PB$, 根据三垂线定理可知 ${AD} \bot  {AB}$ .
        
        因为 $B{C}^{2} + A{B}^{2} = A{C}^{2}$ ,所以 ${BC} \bot  {AB}$ ,根据平面知识可知 ${AD}//{BC}$，又 ${AD} \not \subset$ 平面 ${PBC},{BC} \subset$ 平面 ${PBC}$ ,所以 ${AD}//$ 平面 ${PBC}$ .
    \end{jieda}
\end{enumerate}
\subsection*{证明线面垂直}
\begin{enumerate}
    \item 如图，在$P - {ABCD}$ 中,侧棱 ${PD} \bot$ 长方形底面 ${ABCD}$ , 且 ${PD} = {CD}$ ,过棱 ${PC}$ 的中点 $E$ ,作 ${EF} \bot  {PB}$ 交 ${PB}$ 于点 $F$ ,连接 ${DE},{DF},{BD},{BE}$ . 证明: ${PB} \bot$ 平面 ${DEF}$ .
    
    \begin{center}
    \includegraphics[width=0.2\textwidth]{bodyxb/src/bo_d4ofiljef24c73bbsnjg_0_240_718_330_301_0.jpg}
    \end{center}

    \begin{jieda}
        因为 ${PD} \bot$ 底面 ${ABCD}$ , ${BC} \subset$ 底面 ${ABCD}$ ,所以 ${PD} \bot  {BC}$
        
        由底面 ${ABCD}$ 为长方形,有 ${BC} \bot  {CD}$ ,而 ${PD} \cap  {CD} = D,{PD},{CD} \subset$ 平面 ${PDC}$，所以 ${BC} \bot$ 平面 ${PCD}$ ,而 ${DE} \subset$ 平面 ${PCD}$，所以 ${BC} \bot  {DE}$
        
        又因为 ${PD} = {CD}$ ,点 $E$ 是 ${PC}$ 的中点,所以 ${DE} \bot  {PC}$
        
        而 ${PC} \cap  {BC} = C,{PC},{BC} \subset$ 面 ${PBC}$ ,所以 ${DE} \bot$ 平面 ${PBC}$
        
        而 ${PB} \subset$ 平面 ${PBC}$ ,所以 ${PB} \bot  {DE}$
        
        又 ${PB} \bot  {EF},{DE} \cap  {EF} = E,{DE},{EF} \subset$ 面 ${DEF}$ ,所以 ${PB} \bot$ 平面 ${DEF}$ .
    \end{jieda}

    \item 已知图所示,在四棱锥 $P - {ABCD}$ 中,底面 ${ABCD}$ 是矩形, ${PA} \bot$ 平面 ${ABCD},M,N$ 分别是 ${AB},{PC}$ 的中点, ${PA} = {AD}$ .
    
    \begin{center}
    \includegraphics[width=0.2\textwidth]{bodyxb/src/bo_d4ofiljef24c73bbsnjg_0_240_1804_437_338_0.jpg}
    \end{center}
    \begin{enumerate}
        \item 求证: ${MN}//$ 平面 ${PAD}$
        \item 求证: ${MN}\bot$ 平面 ${PCD}$
    \end{enumerate}
    \begin{jieda}
        \begin{enumerate}
            \item 取$PD$中点$Q$，连接$NQ$，$NQ \overset{/\!/}{=\!\!=} \dfrac{1}{2}CD \overset{/\!/}{=\!\!=} AM$，所以$AMNQ$为平行四边形，故$MN // AQ$
            
            $AQ \subset$平面 ${PAD}$，$MN \not \subset$平面 ${PAD}$，故${MN}//$ 平面 ${PAD}$ 
            \item 因为${PA} = {AD}$，所以$AQ \bot PD$，又因为$MN // AQ$，所以$MN \bot PD$

            因为${PA} \bot$ 平面 ${ABCD}$，所以$PA \bot CD$，又因为$CD \bot AD$，且$PA \cap AD = A$，所以$CD \bot $面$PAD$

            $AQ \subset $面$PAD$，所以$AQ \bot CD$，即$MN \bot CD$

            因为$CD \cap PD = D$，所以${MN} \bot$ 平面 ${PCD}$ .
        \end{enumerate}
    \end{jieda}
    \item 如图,已知多面体 ${ABC}{A}_{1}{B}_{1}{C}_{1}$中${A}_{1}A,{B}_{1}B,{C}_{1}C$ 均垂直于平面 ${ABC}$ ,$\angle {ABC} = {120}^{ \circ  },{A}_{1}A = 4,{C}_{1}C = 1,{AB} = {BC} = {B}_{1}B = 2$ . 证明: $A{B}_{1} \bot$ 平面 ${A}_{1}{B}_{1}{C}_{1}$ .
    
    \begin{center}
    \includegraphics[width=0.2\textwidth]{bodyxb/src/bo_d4ofiljef24c73bbsnjg_2_247_1242_311_382_0.jpg}
    \end{center}
    \begin{jieda}
        根据线面垂直的性质定理可知$AA_1//BB_1//CC_1$，故在直角梯形$ABB_1A_1$中可以算得$A_1B_1 = 2\sqrt{2}, AB_1 = 2\sqrt{2}$，同理，在直角梯形$BCC_1B_1$，中可以算得$B_1C_1 = \sqrt{5}$

        在三角形$ABC$中，可以算得$AC = 2\sqrt{3}$，故在直三角形$ACC_1$中可以算得$AC_1 = \sqrt{13}$

        因此根据勾股定理可知$\angle{AB_1A_1} = \angle{AB_1C_1} = \ang{90}$，即$AB_1 \bot A_1B_1, AB_1 \bot B_1C_1$，所以 $A{B}_{1} \bot$ 平面 ${A}_{1}{B}_{1}{C}_{1}$ .
    \end{jieda}
    \item 如图所示，已知四面体 $A - {BCD}$ 的棱 ${BC} = {AC},{AD} = {BD}$ ,作 ${BE} \bot  {CD}$ 于点 $E$ ,作 ${AH} \bot  {BE}$ 于点 $H$ ,求证: ${AH} \bot$ 平面 ${BCD}$ .
        
    \begin{center}
    \includegraphics[width=0.2\textwidth]{bodyxb/src/bo_d4ogc4bef24c73bbtdng_0_539_853_554_600_0.jpg}
    \end{center}
    
    \begin{jieda}
        取$AB$中点$M$，连接$CM, DM$，因为${BC} = {AC},{AD} = {BD}$，所以$CM \bot AB$，$DM \bot AB$，因此$AB \bot $平面$MCD$，故$AB \bot CD$

        又因为$BE \bot CD$，$AB \cap BE = B$，因此$CD \bot$平面$ABE$，故$CD \bot AH$

        又因为$AH \bot BE$，因此$AH \bot$平面$BCD$.
    \end{jieda}
\end{enumerate}
\subsection*{证明面面平行}
\subsection*{证明面面垂直}

\section*{动点轨迹}
\subsection*{动点与平行结合}
\begin{enumerate}

    \item 在棱长为 2 的正方体 ${ABCD} - {A}_{1}{B}_{1}{C}_{1}{D}_{1}$ 中，点 $E,F$ 分别是棱 ${BC},{C{C}_{1}}$ 的中点，动点 $P$ 是侧面 ${BC}{C}_{1}{B}_{1}$ 内一点，若 ${A}_{1}P//$ 平面 ${AEF}$ ，则点 $P$ 的轨迹长度为\blkda{$\sqrt{2}$}.

    \begin{jieda}
        如下图所示,分别取棱 $B{B}_{1},{B}_{1}{C}_{1}$ 的中点 $M\text{ 、 }N$ ,连接 ${MN},B{C}_{1}$
        
        \begin{center}
        \includegraphics[width=0.2\textwidth]{bodyxb/src/bo_d50clvjef24c73bfi5dg_3_241_1779_393_349_0.jpg}
        \end{center}

        $MN // EF$，$A_1N // AE$，故平面$A_1MN // $平面$AEF$，故$P$的轨迹为$MN$，因此其长度为$\sqrt{2}$.
    \end{jieda}

    \item 已知棱长为 1 的正方体 ${ABCD} - {A}_{1}{B}_{1}{C}_{1}{D}_{1}$ ， $M$ 是 $B{B}_{1}$ 的中点，动点 $P$ 在正方体内部或表面上,且 ${MP}//$ 平面 ${AB}{D}_{1}$ ,则动点 $P$ 的轨迹所形成区域的面积是\blkda{$\dfrac{\sqrt{2}}{2}$}.
    
    \begin{jieda}
        过点 $M$ 作平面 ${AB}{D}_{1}$ 的平行截面,再求解四边形的面积即可.

        \begin{center}
        \includegraphics[width=0.2\textwidth]{bodyxb/src/bo_d50clvjef24c73bfi5dg_5_243_199_492_470_0.jpg}
        \end{center}
    \end{jieda}
    
    \item 如图,在棱长为 2 的正方体 ${ABCD} - {A}_{1}{B}_{1}{C}_{1}{D}_{1}$ 中, $E,F,O$ 分别为正方形 ${AB}{B}_{1}{A}_{1}$，${A}_{1}{B}_{1}{C}_{1}{D}_{1}$ , ${ABCD}$ 的中心,点 $P$ 在正方形 ${ABCD}$ 内 (含边界) 运动,若直线 ${B}_{1}P$ 与平面 ${DEF}$ 无交点, 则点 $P$ 所形成的轨迹长度为\blkda{$\dfrac{2\sqrt{10}}{3}$}.
    
    \begin{center}
    \includegraphics[width=0.2\textwidth]{bodyxb/src/bo_d50clvjef24c73bfi5dg_0_244_716_446_394_0.jpg}
    \end{center}

    \begin{jieda}
        ${B}_{1}P$ 与平面 ${DEF}$ 无交点即${B}_{1}P$ 与平面 ${DEF}$ 平行，故动直线$B_1P$在一与平面 ${DEF}$ 平行的平面内

        $B_1O // DF$，取$CC_1D_1D$的中心$O'$，则$B_1O' // ED$，故平面$B_1OO' // $平面 ${DEF}$
        
        平面$B_1O'O$和平面$ABCD$的交线即为$P$的轨迹，易知$B_1O'$延长后与$AD$相交，设交点为$G$，设直线$OG$和$AB, CD$分别交于$M,N$，则$MN$即为$P$的轨迹，根据勾股定理可以求得其长度为$\dfrac{2\sqrt{10}}{3}$
    \end{jieda}
    
    \item 如图，在四棱锥 $P - {ABCD}$ 中， ${PA}\bot$ 平面 ${ABCD}$ ，底面 ${ABCD}$ 为正方形， ${PA} = {AB} = 4$ ， 点 $E,F$ 分别为 ${CD},{CP}$ 的中点,点 $T$ 为 $\bigtriangleup {PAB}$ 内的一个动点 (包括边界)， ${CT}//$ 平面 ${AEF}$ ，则点 $T$ 的轨迹的长度为\blkda{$\dfrac{2\sqrt{5}}{3}$}.
    
    \begin{center}
    \includegraphics[width=0.2\textwidth]{bodyxb/src/bo_d50clvjef24c73bfi5dg_2_241_611_346_340_0.jpg}
    \end{center}

    \begin{jieda}
        取$AB$中点$M$，则$CM // AE$，连接$BD$交$CM$于$N$，则$\dfrac{BN}{DN} = \dfrac{1}{2}$，取$BP$上靠近$B$的三等分点$Q$，则有$QN // PD // EF$，因此平面$CMQ //$平面$AEF$，故$T$的轨迹为$MQ$，根据余弦定理可求得其长度为$\dfrac{2\sqrt{5}}{3}$
    \end{jieda}

    \item 如图，正三棱柱 ${ABC} - {A}_{1}{B}_{1}{C}_{1}$ 的底面边长是 4，侧棱长是 ${2\sqrt{3}},M$ 为 ${A}_{1}{C}_{1}$ 的中点， $N$ 是侧面 ${BC}{C}_{1}{B}_{1}$ 上一点，且 ${MN}$ //平面 ${AB}{C}_{1}$ ，则线段 ${MN}$ 的最大值为\blkda{$\sqrt{15}$}.

    \begin{center}
    \includegraphics[width=0.2\textwidth]{bodyxb/src/1_93_144_296_355_0.jpg}
    \end{center}

    \begin{jieda}
        取 ${B}_{1}{C}_{1}$ 的中点 $D$ ,取 $B{B}_{1}$ 的中点 $E$ ,连接 ${MD},{DE},{ME},M{B}_{1}$ ,所以 ${DE}//B{C}_{1}$ ,

        又 ${DE} \not \subset$ 面 ${AB}{C}_{1},B{C}_{1} \subset$ 面 ${AB}{C}_{1}$ ,所以 ${DE}//$ 平面 ${AB}{C}_{1}$ ,

        又 $M$ 为 ${A}_{1}{C}_{1}$ 的中点，所以 ${MD}//{A}_{1}{B}_{1}//{AB}$ ，

        又 ${MD} \not \subset$ 面 ${AB}{C}_{1},{AB} \subset$ 面 ${AB}{C}_{1}$ ，所以 ${MD}//$ 平面 ${AB}{C}_{1}$ ，

        \begin{center}
        \includegraphics[width=0.2\textwidth]{bodyxb/src/1_1158_915_301_359_0.jpg}
        \end{center}
        

        又 ${DE} \cap  {MD} = D,{DE} \subset$ 面 ${DEM},{MD} \subset$ 面 ${DEM}$ ,所以平面 ${DEM}//$ 平面 ${AB}{C}_{1}$ , 又因为 $N$ 是侧面 ${BC}{C}_{1}{B}_{1}$ 上一点,且 ${MN}//$ 平面 ${AB}{C}_{1}$ ,

        所以 $N$ 在线段 ${DE}$ 上,

        因为正三棱柱 ${ABC} - {A}_{1}{B}_{1}{C}_{1}$ 的底面边长是 4,侧棱长是 $2\sqrt{3}$

        所以 $B{B}_{1} \bot$ 平面 ${ABC}$ ,因为 $M{B}_{1} \subset$ 平面 ${A}_{1}{B}_{1}{C}_{1}$ ,所以 $B{B}_{1} \bot  M{B}_{1}$

        又 $M$ 为 ${A}_{1}{C}_{1}$ 的中点,所以 ${A}_{1}{C}_{1} \bot  M{B}_{1}$

        所以 $M{B}_{1} = \sqrt{{A}_{1}{B}_{1}^{2} - {A}_{1}{M}^{2}} = \sqrt{{4}^{2} - {2}^{2}} = 2\sqrt{3}$

        则 ${ME} = \sqrt{M{B}_{1}^{2} + {B}_{1}{E}^{2}} = \sqrt{{\left( 2\sqrt{3}\right) }^{2} + {\left( \sqrt{3}\right) }^{2}} = \sqrt{15}$ ,又 ${MD} = 1$

        所以线段 ${MN}$ 的最大值为 $\sqrt{15}$.
    \end{jieda}

    
    \item 如图，在棱长为 2 的正方体 ${ABCD} - {A}_{1}{B}_{1}{C}_{1}{D}_{1}$ 中，点 $E$ 、 $F$ 分别是棱 ${BC}$ ， $C{C}_{1}$ 的中点， $P$ 是侧面 ${BC}{C}_{1}{B}_{1}$ 内(不含边界)一点，若 ${A}_{1}P//$ 平面 ${AEF}$ ，则线段 ${A}_{1}P$ 长度的最小值是\blkda{$\dfrac{3\sqrt{2}}{2}$}.

    \begin{center}
    \includegraphics[width=0.2\textwidth]{bodyxb/src/2_93_1653_340_331_0.jpg}
    \end{center}

    \begin{jieda}
        分别取棱 $B{B}_{1},{B}_{1}{C}_{1}$ 的中点 $M\text{ 、 }N$ ,连接 ${MN},B{C}_{1}$ ,

        $\because M,N,E,F$ 分别为所在棱的中点,则 ${MN}//{B{C}_{1}},{EF}//{B{C}_{1}}$ ,

        $\therefore {MN}//{EF}$ ,又 ${MN} \not \subset$ 平面 ${AEF},{EF} \subset$ 平面 ${AEF}$ ,

        $\therefore {MN}//$ 平面 ${AEF}$ .

        $\because A{A}_{1}//{NE},A{A}_{1} = {NE}$ ， $\therefore$ 四边形 ${AEN}{A}_{1}$ 为平行四边形，

        $\therefore {A}_{1}N//{AE}$ ,

        又 ${A}_{1}N \not \subset$ 平面 ${AEF},{AE} \subset$ 平面 ${AEF}$ ,

        \begin{center}
        \includegraphics[width=0.2\textwidth]{bodyxb/src/3_997_513_392_375_0.jpg}
        \end{center}
        

        $\therefore {A}_{1}N//$ 平面 ${AEF}$ ,又 ${A}_{1}N \cap  {MN} = N$ ,

        $\therefore$ 平面 ${A}_{1}{MN}//$ 平面 ${AEF}$ .

        $\because P$ 是侧面 ${BC}{C}_{1}{B}_{1}$ 内一点,且 ${A}_{1}P//$ 平面 ${AEF}$ ,

        $\therefore$ 点 $P$ 必在线段 ${MN}$ 上.

        在 $R{t}_{ \bigtriangleup  }{A}_{1}{B}_{1}M$ 中， ${A}_{1}M = \sqrt{{A}_{1}{B}_{1}^{2} + {B}_{1}{M}^{2}} = \sqrt{{2}^{2} + {1}^{2}} = \sqrt{5}$ .

        同理,在 ${Rt} \bigtriangleup  {A}_{1}{B}_{1}N$ 中,可得 ${A}_{1}N = \sqrt{5}$ ,

        $\therefore \bigtriangleup {A}_{1}{MN}$ 为等腰三角形.

        当点 $P$ 为 ${MN}$ 中点 $O$ 时,即 ${A}_{1}P \bot  {MN}$ ,此时 ${A}_{1}P$ 最短;

        又 ${A}_{1}O = \sqrt{{A}_{1}{M}^{2} - O{M}^{2}} = \sqrt{{\left( \sqrt{5}\right) }^{2} - {\left( \dfrac{\sqrt{2}}{2}\right) }^{2}} = \dfrac{3\sqrt{2}}{2}$
    \end{jieda}
\end{enumerate}
\subsection*{动点与垂直结合}
\begin{enumerate}
    \item 如图，正四棱锥 $P - {ABCD}$ 的底面边长为 2,高为 2, $E$ 是边 ${BC}$ 的中点，动点 $Q$ 在表面上运动， 并且总保持 ${QE}\bot {AC}$ ,则动点 $Q$ 的轨迹的周长为\blkda{$\sqrt{2} + \sqrt{6}$}.
    \begin{center}
        \includegraphics[width=0.2\linewidth]{bodyxb//src/dongdian1.png}
    \end{center}
    \begin{jieda}
        $Q$的轨迹为过$E$的$AC$的垂面在正四棱锥 $P - {ABCD}$上的截面

        取$PC$和$CD$的中点$F,G$，取$EG$中点$M$，则$FM // PO$，故$FM \bot $平面$ABCD$，故$FM  \bot AC$，且$EG \bot AC$，所以$AC \bot $平面$EFG$，故$P$的轨迹为$\triangle EFG$，其周长为$\sqrt{2}+\sqrt{6}$
    \end{jieda}
    
    \item 已知四棱柱 ${ABCD} - {A}_{1}{B}_{1}{C}_{1}{D}_{1}$ 的底面为菱形， $A{A}_{1} \bot$ 底面 ${ABCD},A{A}_{1} = 6,{AB} = 6$ ， $\angle {BCD} = {60}^{ \circ  }$ ,点 $M$ 是线段 ${BC}$ 上靠近 $C$ 的三等分点,点 $N$ 在四棱柱 ${ABCD} - {A}_{1}{B}_{1}{C}_{1}{D}_{1}$ 的表面上移动，且 ${MN}\bot {A}_{1}C$ ，则动点 $N$ 的轨迹长度为\blkda{$2 + 2\sqrt{13}$}.

    \begin{jieda}
        如图所示
            
        \begin{center}
        \includegraphics[width=0.2\textwidth]{bodyxb/src/bo_d50clvjef24c73bfi5dg_8_241_1196_358_343_0.jpg}
        \end{center}
    
        取$CD$上靠近$C$的三等分点，则$MF \bot AC$，取$CC_1$中点$G$，故平面$MFG \bot $平面$ACC_1A_1$，因为$AC \bot OG$，所以$AC \bot$平面$GMF$，故$N$的轨迹为$\triangle MFG$，则动点 $N$ 的轨迹长度为 ${MF} + {MG} + {FG} = 2 + 2\sqrt{13}$ .
    \end{jieda}
\end{enumerate}
\subsection*{动点与距离结合}
\begin{enumerate}
    \item 已知正三棱锥 $P - {ABC}$ 的六条棱长均为 $6,S$ 是 $\bigtriangleup {ABC}$ 及其内部的点构成的集合. 设集合 $T = \left\{  {Q \mid  {PQ} \leq  2\sqrt{7},Q \in  S}\right\}$ ,则 $T$ 表示的区域的面积为\blkda{${2\pi } + 3\sqrt{3}$}
    \begin{jieda}
        设顶点 $P$ 在底面上的投影为 $O$ ,连接 ${BO}$ ,则 $O$ 为三角形 ${ABC}$ 的中心,
        
        且 ${BO} = \dfrac{2}{3} \times  6 \times  \dfrac{\sqrt{3}}{2} = 2\sqrt{3}$ ,故 ${PO} = \sqrt{P{B}^{2} - B{O}^{2}} = \sqrt{{6}^{2} - {\left( 2\sqrt{3}\right) }^{2}} = 2\sqrt{6}$ .
        
        \begin{center}
        \includegraphics[width=0.2\textwidth]{bodyxb/src/bo_d50clvjef24c73bfi5dg_10_241_286_286_391_0.jpg}
        \end{center}
        
        
        因为 ${PQ} \leq  2\sqrt{7}$ ,故 ${OQ} = \sqrt{P{Q}^{2} - P{O}^{2}} \leq  \sqrt{{\left( 2\sqrt{7}\right) }^{2} - {\left( 2\sqrt{6}\right) }^{2}} = 2$ ,
        
        故 $Q$ 在以 $O$ 为圆心,2 为半径的圆及其内部,
        
        而三角形 ${ABC}$ 内切圆的圆心为 $O$ ,半径为 $\dfrac{1}{3} \times  6 \times  \dfrac{\sqrt{3}}{2} = \sqrt{3}$ ,而 $\sqrt{3} < 2 < 2\sqrt{3}$ ,
        
        故 $T$ 表示的区域为以 $O$ 为圆心,2 为半径的圆在 $\bigtriangleup  {ABC}$ 内(包含边界)的部分，
        
        设该圆与 ${AB}$ 交于 $D,E$ 两点,则 ${DE} = 2\sqrt{{2}^{2} - {\left( \sqrt{3}\right) }^{2}} = 2$ ,
        
        即 $\bigtriangleup  {ODE}$ 为正三角形,则弧 ${DE}$ 所在的弓形的面积为 $\dfrac{1}{2} \times  \dfrac{\pi }{3} \times  4 - \dfrac{\sqrt{3}}{4} \times  {2}^{2} = \dfrac{2\pi }{3} - \sqrt{3}$ ,
        
        故 $T$ 表示的区域的面积为 $\pi  \times  {2}^{2} - 3\left( {\dfrac{2\pi }{3} - \sqrt{3}}\right)  = {2\pi } + 3\sqrt{3}$
    \end{jieda}
\end{enumerate}
\subsection*{动点与角度结合}
\begin{enumerate}
    \item 已知平面 ${ABCD} \bot$ 平面 ${ADEF}$ ， ${AB} \bot  {AD},{CD} \bot  {AD}$ ，且 ${AB} = 3$ ， ${AD} = {CD} = 6$ ， ${ADEF}$ 是正方形，在正方形 ${ADEF}$ 内部(含边界)有一点 $M$ ，满足 ${MB}$ ， ${MC}$ 与平面 ${ADEF}$ 所成的角相等，则点 $M$ 的轨迹长度为\blkda{$\dfrac{4\pi }{3}$}.

    \begin{jieda}
        因为平面 ${ABCD} \bot$ 平面 ${ADEF}$ ,平面 ${ABCD} \cap$ 平面 ${ADEF} = {AD}$ ,
        
        ${AB} \bot  {AD},{CD} \bot  {AD}$ ,
        
        且 ${AB},{CD} \subset$ 平面 ${ABCD}$ ,所以 ${AB} \bot$ 平面 ${ADEF},{CD} \bot$ 平面 ${ADEF}$ ,
        
        可知 $\angle {BMA}$ 和 $\angle {CMD}$ 分别是直线 ${MB},{MC}$ 与平面 ${ADEF}$ 所成的角,
        
        则 $\angle {BMA} = \angle {CMD}$ ,即 $\tan \angle {BMA} = \tan \angle {CMD}$ ,
        
        可得 $\dfrac{AB}{AM} = \dfrac{CD}{MD}$ ,所以 ${MD} = {2AM}$ .

        故$M$的轨迹为阿氏圆（的一部分）

        以 $A$ 为原点建立平面直角坐标系如下图所示,
        \begin{center}
        \includegraphics[width=0.2\textwidth]{bodyxb/src/bo_d50clvjef24c73bfi5dg_9_569_488_343_285_0.jpg}
        \end{center}

        由 ${MD} = {2AM}$ 得 $\sqrt{{\left( x - 6\right) }^{2} + {y}^{2}} = 2\sqrt{{x}^{2} + {y}^{2}}$ ,化简得 ${\left( x + 2\right) }^{2} + {y}^{2} = {16}$
        
        由于点 $M$ 在第一象限内，所以 $M$ 点的轨迹是以 $G\left( {-2,0}\right)$ 为圆心，半径为 4 的圆在第一象限的部分,
        
        令 $x = 0$ 代入原方程,解得 $y =  \pm  2\sqrt{3}$ ,故 $H\left( {0,2\sqrt{3}}\right)$ ,
        
        由于 ${GA} = 2$ ，所以 $\angle {HGA} = \dfrac{\pi }{3}$ ，
        
        所以点 $M$ 的轨迹长度为 $\dfrac{\pi }{3} \times  4 = \dfrac{4\pi }{3}$ .
    \end{jieda}
\end{enumerate}

\subsection*{截面面积}
\begin{enumerate}
    \item 如图,正方体的棱长是 $a$ ,点 $E\text{ 、 }F$ 分别是两条棱的中点. 求四边形 ${BDEF}$ 的面积.
    
    \begin{center}
    \includegraphics[width=0.2\textwidth]{bodyxb/src/bo_d4lnl5ref24c739do320_11_241_208_336_326_0.jpg}
    \end{center}

    \begin{jieda}
        易知四边形 ${BDEF}$ (图中阴影部分)是一个梯形.
    
        因为正方体的棱长为 $a$ ，所以 ${BD} = \sqrt{2}a,{EF} = \dfrac{\sqrt{2}}{2}a,{ED} = {BF} = \dfrac{\sqrt{5}}{2}a$ ， 如图所示， ${DM} = \dfrac{1}{2}\left( {{BD} - {EF}}\right)  = \dfrac{\sqrt{2}}{4}a$ ，
        
        \begin{center}
        \includegraphics[width=0.2\textwidth]{bodyxb/src/bo_d4lnl5ref24c739do320_11_243_1855_408_242_0.jpg}
        \end{center}
        
        
        而梯形 ${BDEF}$ 的高 ${EM} = \sqrt{E{D}^{2} - D{M}^{2}} = \sqrt{\dfrac{5}{4}{a}^{2} - \dfrac{1}{8}{a}^{2}} = \dfrac{3\sqrt{2}}{4}a$ ，
        
        可得梯形 ${BDEF}$ 的面积为 $S = \dfrac{1}{2} \times  \left( {\dfrac{\sqrt{2}}{2}a + \sqrt{2}a}\right)  \times  \dfrac{3\sqrt{2}}{4}a = \dfrac{9}{8}{a}^{2}$ .
    \end{jieda}

    \item 如图，在四棱锥 $P - {ABCD}$ 中，底面 ${ABCD}$ 是正方形. ${PD} \bot$ 平面 ${ABCD}$ ， ${PD} = {AD} = 2$ ，点 $E$ 是棱 ${PC}$ 的中点,平面 ${ABE}$ 与棱 ${PD}$ 交于点 $F$，求四边形$ABEF$的面积.
    
    \begin{center}
    \includegraphics[width=0.2\textwidth]{bodyxb/src/bo_d4lnl5ref24c739do320_13_241_198_342_344_0.jpg}
    \end{center}
    \begin{jieda}
        ${CD}//{AB},{AB} \subset$ 平面 ${ABE},{CD} \subset$ 平面 ${ABE}$，所以${CD}//$ 平面 ${ABE}$
        
        又 ${CD} \subset$ 平面 ${PCD}$ ,平面 ${PCD} \cap$ 平面 ${ABE} = {EF}$，故${CD}//{EF}$，故$EF$为$\triangle PAC$的中位线，所以$EF = 1$
        
        由于$CD//AB$，根据平行的传递性可知${AB}//{EF}$，故$ABEF$为梯形

        根据三垂线定理可知$AB \bot AF$，故$ABEF$为直角梯形，$EF = 1, AB = 2, AF = \sqrt{5}$，故面积为$\dfrac{3\sqrt{5}}{2}$.
    \end{jieda}
\end{enumerate}

\section*{角的最值}
\begin{enumerate}
    \item 已知正三棱柱 ${ABC} - {A}^{\prime }{B}^{\prime }{C}^{\prime }$ 的所有棱长均为 2,且点 $P$ 在 ${B}^{\prime }{C}^{\prime }$ 上运动,则直线 ${AP}$ 与平面${ABC}$ 所成角的最大正弦值为\blkda{$\dfrac{2\sqrt{7}}{7}$}
    \begin{jieda}
        根据最大角定理可知且点 $P$ 在 ${B}^{\prime }{C}^{\prime }$ 上运动,则直线 ${AP}$ 与平面${ABC}$ 所成角的最大值为平面$B'C'A$和平面$BCA$所成锐二面角的的大小.
    \end{jieda}
    \item 如图,在四边形 ${ABCD}$ 中, ${AB} = {BC} = {CD} = 2,{AD} = \sqrt{5}$ ,对角线 ${BD} = 3,E$ 是直线 ${CD}$ 上任一点，将 $\bigtriangleup {ABD}$ 沿 ${BD}$ 翻折成 $\bigtriangleup {A}^{\prime }{BD}$ ,使二面角 ${A}^{\prime } - {BD} - C$ 的大小为 ${120}^{ \circ  }$ , 设异面直线 ${A}^{\prime }D$ 和 ${BE}$ 所成的角为 $\alpha$ ,则 $\sin \alpha$ 的最小值是\blkda{$\dfrac{\sqrt{3}}{3}$}.

    \twotu{\includegraphics[width=0.5\textwidth]{bodyxb/src/bo_d1ig044601uc738hj6v0_69_813_1076_298_264_0.jpg}}{\includegraphics[width=0.5\textwidth]{bodyxb/src/bo_d1ig044601uc738hj6v0_69_1138_1072_342_264_0.jpg}}
    \begin{jieda}
        当$sin\alpha$最小的时候，$\alpha$最小，根据最小角定理可知此时${A}^{\prime }D$ 和 ${BE}$ 所成的角等于${A}^{\prime }D$ 和平面$BCD$所成角，设$A'$在面$BCD$内射影为$H$，故计算可得$sin\alpha = \dfrac{A'H}{A'D} = \dfrac{\dfrac{\sqrt{15}}{3}}{\sqrt{5}} = \dfrac{\sqrt{3}}{3}$
    \end{jieda}
    \item 在四面体 ${ABCD}$ 中, $\bigtriangleup {BCD}$ 为等边三角形, $\angle {ADB} = \dfrac{\pi }{2}$ ,二面角 $B - {AD} - C$ 的大小为 $\alpha$ , 则 $\alpha$ 的取值范围是\blkda{$\left(0, \dfrac{\pi}{3}\right]$}
    \begin{center}
    \includegraphics[width=0.2\textwidth]{bodyxb/src/bo_d1ig044601uc738hj6v0_69_456_1053_312_290_0.jpg}
    \end{center}
    \begin{jieda}
        因为$DB \bot AD$，所以$DB$和平面$ACD$的夹角即为$B - {AD} - C$的平面角，根据最小角定理可知$\alpha \leq \angle{BDC} = \dfrac{\pi}{3}$
    \end{jieda}
    
    \item 如图,在三棱柱 ${ABC} - {A}_{1}{B}_{1}{C}_{1}$ 中, ${AB}\text{ 、 }{AC}\text{ 、 }A{A}_{1}$ 两两互相垂直, ${AB} = {AC} = A{A}_{1},M\text{ 、 }N$ 是线段 $B{B}_{1}\text{ 、 }C{C}_{1}$ 上的点,平面 ${AMN}$ 与平面 ${ABC}$ 所成(锐)二面角为 $\dfrac{\pi }{6}$ . 当 $\left| {{B}_{1}M}\right|$ 最小时, $\angle {AMB} =$ \blkda{$\dfrac{\pi }{3}$}
    
    \begin{center}
    \includegraphics[width=0.2\textwidth]{bodyxb/src/bo_d1ig044601uc738hj6v0_69_172_1006_243_337_0.jpg}
    \end{center}
    \begin{jieda}
        易知三棱柱 ${ABC} - {A}_{1}{B}_{1}{C}_{1}$为直三棱柱，$\left| {{B}_{1}M}\right|$ 最小即$\angle{AMB}$最小.
        
        根据最大角定理可知$AM$和平面$ABC$的线面角（$MAB$）的最大值为平面 ${AMN}$ 与平面 ${ABC}$ 所成(锐)二面角，故$\angle{MAB} \leq \dfrac{\pi}{6}$，即$\angle{AMB} \geq \dfrac{\pi}{3}$
    \end{jieda}
    \item 设三棱锥 $V - {ABC}$ 的底面是正三角形,侧棱长均相等, $P$ 是棱 ${VA}$ 上的点 (不含端点). 记直线 ${PB}$ 与直线 ${AC}$ 所成的角为 $\alpha$ ,直线 ${PB}$ 与平面 ${ABC}$ 所成的角为 $\beta$ ,二面角 $P - {AC} - B$ 的平面角为 $\gamma$ ,则 \dotfill{(\kh{B})}
    \xx{$\beta  < \gamma ,\alpha  < \gamma$}{$\beta  < \alpha ,\beta  < \gamma$}{$\beta  < \alpha ,\gamma  < \alpha$}{$\alpha  < \beta ,\gamma  < \beta$}
    \begin{jieda}
        根据最小角定理可知$\alpha \geq \beta$，易知不满足取等条件，故$\beta < \alpha$

        \ps 因为$PB$并不在二面角$P - {AC} - B$的半平面内，因此不适用最大角定理

        过$P$做$PM \bot $面$ABC$，做$MN \bot AC$，则$\beta = \angle{PBM}$，$\gamma = \angle{PNM}$，故$tan\gamma = \dfrac{PM}{MN}$，$tan\beta = \dfrac{PM}{BM}$

        $V - {ABC}$为正四面体，根据正四面体的性质可知$M$在$\angle{BAC}$的中垂线上，有$MB = MC$，因为$MC > MN$，所以$tan\gamma > tan\beta$
    \end{jieda}
    
    \item 如图,在 $\bigtriangleup {ABC}$ 中, ${AB} = 1,{BC} = 2\sqrt{2},\angle B = \dfrac{\pi }{4}$ ,将 $\bigtriangleup {ABC}$ 绕边 ${AB}$ 翻转至 $\bigtriangleup {ABP}$ ,使平面 ${ABP} \bot$ 平面 ${ABC},D$ 是边 ${BC}$ 的中点,设 $Q$ 是线段 ${PA}$ 上的动点,当 ${PC}$ 与 ${DQ}$ 所成的角取得最小值时,线段 ${AQ}$ 的长度为\blkda{$\dfrac{2\sqrt{5}}{5}$}.
    
    \begin{center}
    \includegraphics[width=0.2\textwidth]{bodyxb/src/bo_d1ig044601uc738hj6v0_69_1169_124_331_299_0.jpg}
    \end{center}
    \begin{jieda}
        根据最小角定理可知，当 ${PC}$ 与 ${DQ}$ 所成的角取得最小值时，$DQ$与$PC$在面$APD$内射影平行

        延长$BA$至$O$使得$AO = 1$，则$\angle{BOC} = \dfrac{\pi}{2}$，故$CO \bot $平面$BOP$，且$PO \bot $平面$BOC$（$O$为墙角）

        将图形补成立方体，补全$PAD$对应截面，并做出$PC$在面内的投影，根据比例关系可以求得$AQ = \dfrac{2\sqrt{5}}{5}$
    \end{jieda}
    
\end{enumerate}

\section*{综合}
\begin{enumerate}
    \item 如图,在棱长为 1 的正方体 ${ABCD} - {A}_{1}{B}_{1}{C}_{1}{D}_{1}$ 中,点 $P$ 满足 $\vv{BP} = \lambda \vv{BC} + \mu \vv{B{B}_{1}},\lambda  \in  \left\lbrack  {0,1}\right\rbrack  ,\mu  \in  \left\lbrack  {0,1}\right\rbrack$ ,则下列说法正确的是\blkda{\ding{172} \ding{173}}
    
    \begin{center}
    \includegraphics[width=0.2\textwidth]{bodyxb/src/bo_d4q3m9bef24c73bcm6mg_0_241_633_357_342_0.jpg}
    \end{center}

    \begin{dingautolist}{172}
        \item 当 $\lambda  = \mu  \neq  0$ 时， ${BP}//$ 平面 $A{B}_{1}{D}_{1}$ 
        
        \item 当 $\lambda  = \mu$ 时， ${AP}\bot {B}_{1}C$ 
        
        \item 当 $\lambda  + \mu  = 1$ 时， ${AP}$ 长度的最小值为 $\dfrac{\sqrt{3}}{2}$ 
        
        \item 当 $\lambda  + \mu  = 1$ 时,存在点 $P$ ,使得 ${AP}$ 与平面 ${BC}{C}_{1}{B}_{1}$ 所成的角为 $\dfrac{\pi }{3}$ .
    \end{dingautolist}
    \begin{jieda}
        \begin{dingautolist}{172}
            \item 当 $\lambda  = \mu  \neq  0$ 时， $P$ 为线段 $B{C}_{1}$ 上的点(不含 $B$ 点)，
            
            显然 ${BP}//A{D}_{1}$ ,又 ${BP} \not \subset$ 平面 $A{B}_{1}{D}_{1},A{D}_{1} \subset$ 平面 $A{B}_{1}{D}_{1}$，所以 ${BP}//$ 平面 $A{B}_{1}{D}_{1}$

            \item 当 $\lambda  = \mu$ 时， $P$ 为线段 $B{C}_{1}$ 上的点，又 ${AB}$ 垂直右侧面，
            
            所以 ${AP}$ 在右侧面的射影为 $B{C}_{1}$ ,又 $B{C}_{1} \bot  {B}_{1}C$ , 根据三垂线定理可知 ${AP} \subset$ 平面 ${AB}{C}_{1}$，所以 ${AP}\bot {B}_{1}C$

            \item 当 $\lambda  + \mu  = 1$ 时， $P$ 为线段 ${B}_{1}C$ 上的点，又三角形 $A{B}_{1}C$ 是边长为 $\sqrt{2}$ 的等边三角形，所以 ${AP}$ 长度的最小值为该等边三角形的高,即为 $\sqrt{2} \times  \dfrac{\sqrt{3}}{2} = \dfrac{\sqrt{6}}{2}$

            \item 当 $\lambda  + \mu  = 1$ 时， $P$ 为线段 ${B}_{1}C$ 上的点，又 ${AB}$ 垂直右侧面，所以 ${AP}$ 与平面 ${BC}{C}_{1}{B}_{1}$ 所成的角为 $\angle {APB}$

            $\tan \angle {APB} = \dfrac{AB}{BP} = \dfrac{1}{BP}$ ,
        
            所以当 ${BP}$ 取最小值 $\dfrac{1}{2}B{C}_{1} = \dfrac{\sqrt{2}}{2}$ 时, $\tan \angle {APB} = \dfrac{1}{BP}$ 取得最大值 $\sqrt{2} < \tan \dfrac{\pi }{3} = \sqrt{3}$ ,
            
            所以 $\angle {APB}$ 小于 $\dfrac{\pi }{3}$ ,所以不存在点 $P$ ,使得 ${AP}$ 与平面 ${BC}{C}_{1}{B}_{1}$ 所成的角为 $\dfrac{\pi }{3}$
        \end{dingautolist}
        
        \begin{center}
        \includegraphics[width=0.2\textwidth]{bodyxb/src/bo_d4q3m9bef24c73bcm6mg_1_240_1132_357_342_0.jpg}
        \end{center}
    \end{jieda}

    \item 如图，四边形 ${ABCD}$ 是正方形， ${AF} \bot$ 平面 ${ABCD}$ ， ${AF}//{BE},{AB} = {BE} = 2,{AF} = 1$ .

    \begin{enumerate}
        \item 求证:平面 ${ACF} \bot$ 平面 ${BDE}$
        \item 求点 $C$ 到平面 ${BDF}$ 的距离
    \end{enumerate}

    \begin{center}
    \includegraphics[width=0.2\textwidth]{bodyxb/src/2_1041_771_335_343_0.jpg}
    \end{center}
    \begin{jieda}
        \begin{enumerate}
            \item 因为 ${AF} \bot$ 平面 ${ABCD}$ ， ${AC} \subset$ 平面 ${ABCD}$ ，所以 ${AF} \bot  {AC}$ ， 因为 ${AF}//{BE}$ ,所以 ${BE} \bot  {AC}$,

            因为四边形 ${ABCD}$ 是正方形,所以 ${BD} \bot  {AC}$ ,

            因为 ${BD} \cap  {BE} = B,{BD},{BE} \subset$ 平面 ${BDE}$ ,所以 ${AC} \bot$ 平面 ${BDE}$ 因为 ${AC} \subset$ 平面 ${ACF}$ ,所以平面 ${ACF} \bot$ 平面 ${BDE}$.
            \item 过点 $C$ 作 ${CG} \bot  {OF}$ 于点 $G$ ，

            因为 ${AF} \bot  {BD},{BD} \bot  {AC}$ ， ${AF} \cap  {AC} = A$ ， ${AF},{AC} \subset$ 平面 ${ACF}$ ，

            所以 ${BD} \bot$ 平面 ${ACF}$ ,

            因为 ${CG} \subset$ 平面 ${ACF}$ ,所以 ${BD} \bot  {CG}$ ,

            因为 ${CG} \bot  {OF},{BD} \cap  {OF} = O,{BD},{OF} \subset$ 平面 ${BDF}$ ,所以 ${CG} \bot$ 平面 ${BDF}$ ,

            则 ${CG}$ 的长即为点 $C$ 到平面 ${BDF}$ 的距离,

            又 ${AB} = 2,{AF} = 1$ ,则 ${OA} = {OC} = \sqrt{2},{OF} = \sqrt{A{O}^{2} + A{F}^{2}} = \sqrt{1 + 2} = \sqrt{3}$ ,

            则 ${CG} = {OC}\sin \angle {COG} = {OC}\sin \angle {AOF} = \sqrt{2} \times  \dfrac{1}{\sqrt{3}} = \dfrac{\sqrt{6}}{3}$ ，故点 $C$ 到平面 ${BDF}$ 的距离为 $\dfrac{\sqrt{6}}{3}$.

            \begin{center}
            \includegraphics[width=0.2\textwidth]{bodyxb/src/2_260_1806_298_346_0.jpg}
            \end{center}
        \end{enumerate}
    \end{jieda}

    \item 如图，等腰直角三角形 ${ABC}$ 中， $\angle {ACB} = \dfrac{\pi }{2}$ ， $D$ 为 ${AC}$ 中点， $E$ ， $F$ 分别为 ${BA}$ ， ${BC}$ 边上的动点,且 ${EF}//{AC}$ ,将三角形 ${BEF}$ 沿 ${EF}$ 折起,使点 $B$ 折至点 $P$ 的位置,二面角 $P - {EF} - C$ 大小为 $\dfrac{\pi }{3}$ ,连接 ${PA},{PC},{PD}$.
    \begin{enumerate}
        \item 求证: ${AC}\bot$ 平面 ${PCF}$
        \item 若点 $F$ 为 ${BC}$ 中点，求异面直线 ${PD}$ 与 ${CF}$ 所成的角的大小
        \item 试判断直线 ${PE}$ 与平面 ${PAC}$ 所成的角的大小是否存在最大值? 若存在, 求出最大值; 若不存在, 请说明理由.
    \end{enumerate}

    \begin{center}
    \includegraphics[width=0.2\textwidth]{bodyxb/src/0_1059_1703_364_231_0.jpg}
    \end{center}
    \begin{jieda}
        \begin{enumerate}
            \item 由 ${AC} \bot  {BC},{EF}//{AC}$ ,得 ${EF} \bot  {BC}$ ,则 ${EF} \bot  {FC},{EF} \bot  {PF}$ , 因此 ${AC} \bot  {FC},{AC} \bot  {PF}$ ,而 ${PF} \cap  {CF} = F,{PF},{CF} \subset$ 平面 ${PCF}$ , 所以 ${AC} \bot$ 平面 ${PCF}$.
            \item 连接 ${DE}$ ，由点 $F$ 为 ${BC}$ 中点， ${EF}//{AC}$ ，得 $E$ 是 ${AB}$ 的中点，而 $D$ 为 ${AC}$ 中点， 则 ${DE}//{CF},{DE} = {CF},\angle {PDE}$ 是异面直线 ${PD}$ 与 ${CF}$ 所成的角或其补角,

            由(1)得 $\angle {PFC}$ 是二面角 $P - {EF} - C$ 的平面角， $\angle {PFC} = \dfrac{\pi }{3}$ ，令 ${PF} = {CF} = a$ ，

            则 $\bigtriangleup {PCF}$ 是正三角形， ${CD} = {PC} = a$ ，由( 1 )得 ${AC} \bot  {PC}$ ，于是 ${PD} = \sqrt{2}a$ ，

            而 ${DE} = a,{PE} = {BE} = \sqrt{2}a$ ，在等腰 $\bigtriangleup  {PDE}$ 中， $\cos \angle {PDE} = \dfrac{\dfrac{1}{2}{DE}}{PD} = \dfrac{\sqrt{2}}{4}$ ，则 $\angle {PDE} = \arccos \dfrac{\sqrt{2}}{4}$ ， 所以异面直线 ${PD}$ 与 ${CF}$ 所成的角的大小为 $\arccos \dfrac{\sqrt{2}}{4}$ .

            \begin{center}
            \includegraphics[width=0.2\textwidth]{bodyxb/src/1_243_758_359_222_0.jpg}
            \end{center}
            \item 作 ${FG} \bot  {PC}$ 于点 $G$ ,由 ${AC} \bot$ 平面 ${PCF},{GF} \subset$ 平面 ${PCF}$ ,得 ${AC} \bot  {FG}$ ,

            而 ${AC} \cap  {PC} = C,{AC},{PC} \subset$ 平面 ${PAC}$ ,则 ${FG} \bot$ 平面 ${PAC}$ ,

            又 ${EF}//{AC},{EF} \not \subset$ 平面 ${PAC},{AC} \subset$ 平面 ${PAC}$ ,于是 ${EF}//$ 平面 ${PAC}$ ,

            点 $E$ 到平面 ${PAC}$ 距离等于点 $F$ 到平面 ${PAC}$ 距离 ${FG}$ ,

            设 ${BC} = 1,{BF} = x$ ,直线 ${PE}$ 与平面 ${PAC}$ 所成的角为 $\theta$ ,

            则 ${CF} = 1 - x,{PC} = \sqrt{{x}^{2} + {\left( 1 - x\right) }^{2} - {2x}\left( {1 - x}\right) \cos \dfrac{\pi }{3}} = \sqrt{3{x}^{2} - {3x} + 1},{FG} = \dfrac{\sqrt{3}x\left( {1 - x}\right) }{2\sqrt{3{x}^{2} - {3x} + 1}}$ ,

            而 ${PE} = \sqrt{2}{BF} = \sqrt{2}x$ ,因此 $\sin \theta  = \dfrac{FG}{PE} = \dfrac{\sqrt{6}\left( {1 - x}\right) }{4\sqrt{3{x}^{2} - {3x} + 1}}$ ,令 $1 - x = t \in  \left( {0,1}\right)$ ,

            则 $\dfrac{3{x}^{2} - {3x} + 1}{{\left( 1 - x\right) }^{2}} = \dfrac{1}{{t}^{2}} - \dfrac{3}{t} + 3 = {\left( \dfrac{1}{t} - \dfrac{3}{2}\right) }^{2} + \dfrac{3}{4} \geq  \dfrac{3}{4}$ ,当且仅当 $t = \dfrac{2}{3},x = \dfrac{1}{3}$ 时取等号,

            因此 ${\left( \dfrac{1 - x}{\sqrt{3{x}^{2} - {3x} + 1}}\right) }_{\max } = \sqrt{\dfrac{4}{3}} = \dfrac{2\sqrt{3}}{3},{\left( \sin \theta \right) }_{\max } = \dfrac{\sqrt{6}}{4} \cdot  \dfrac{2\sqrt{3}}{3} = \dfrac{\sqrt{2}}{2}$ ,又 $\theta$ 为锐角,则 ${\theta }_{\max } = \dfrac{\pi }{4}$ ,

            所以直线 ${PE}$ 与平面 ${PAC}$ 所成的角的大小存在最大值,最大值为 $\dfrac{\pi }{4}$.
        \end{enumerate}
    \end{jieda}
\end{enumerate}